\documentclass[11pt]{article}
\usepackage[T1]{fontenc}
\usepackage{lmodern}
\usepackage{amsmath,amssymb,amsthm,mathtools}
\usepackage[margin=1in]{geometry}
\usepackage{booktabs,array,longtable}
\usepackage[authoryear,round]{natbib}
\usepackage{microtype}
\usepackage[colorlinks=true,linkcolor=blue,citecolor=blue,urlcolor=blue]{hyperref}
\usepackage{enumitem}
\setlist{itemsep=3pt,topsep=5pt}
\newtheorem{theorem}{Theorem}[section]
\newtheorem{proposition}[theorem]{Proposition}
\newtheorem{lemma}[theorem]{Lemma}
\newtheorem{corollary}[theorem]{Corollary}
\theoremstyle{definition}
\newtheorem{definition}[theorem]{Definition}
\theoremstyle{remark}
\newtheorem{remark}[theorem]{Remark}
\newcommand{\Q}{\mathbb Q}
\newcommand{\R}{\mathbb R}
\newcommand{\C}{\mathbb C}
\newcommand{\Z}{\mathbb Z}
\newcommand{\N}{\mathbb N}
\newcommand{\GL}{\mathrm{GL}}
\newcommand{\Gal}{\operatorname{Gal}}
\newcommand{\ord}{\operatorname{ord}}
\newcommand{\Ind}{\operatorname{Ind}}
\newcommand{\supp}{\operatorname{supp}}
\newcommand{\Gr}{\Gamma_{\R}}
\newcommand{\Gc}{\Gamma_{\C}}
\newcommand{\one}{\mathbf 1}
\newcommand{\rh}{\rho_5}
\newcommand{\dd}{\,d}
\newcommand{\Rea}{\operatorname{Re}}
\newcommand{\Ima}{\operatorname{Im}}
\newcommand{\workbenchurl}{https://github.com/KokunoYumeto/zeta-and-l-functions-workbench}
\newcommand{\recordurl}{https://github.com/KokunoYumeto/mathematics-workbenches}
\numberwithin{equation}{section}
\title{Real zeros and Artin's conjecture for some\texorpdfstring{ $S_5$}{ S5}-extensions\\[3pt]\large after Calegari}
\author{}
\date{}
\begin{document}
\maketitle
\begin{abstract}
Calegari constructs a tempered cuspidal representation on $\GL_5/\Q$ associated on a density-one set of primes to certain five-dimensional $S_5$ Artin representations. Pilloni--Stroh's later modularity theorem extends this construction to every $S_5$-extension with complex conjugation of type $(12)(34)$, without a condition at five. Calegari's final comparison theorem assumes nonvanishing of a specified Dedekind zeta function throughout the real critical interval. We prove a refinement that permits central zeros of arbitrary order and only excludes common off-central real zeros of two explicit quadratic Hecke functions. The first is $B=\zeta_H/\zeta_F$, for the quadratic field $E$, the sextic resolvent $F$, and $H=EF$. The second is $C=\zeta_{J'}/\zeta_J$, where the pair stabilizers $S_2\times S_3$ and $S_2\times A_3$ define fields of degrees ten and twenty. Odd-rank exterior-square holomorphy shows that every comparison pole is a common zero of $B$ and $C$, with its order bounded by each zero order. The proof then turns the pole orders of the completed comparison quotient into an integer-valued positive-definite function on the discrete additive real group. Cohen's idempotent theorem and the structure of discrete coset-ring sets force an arithmetic-progression description of its divisor. Hadamard factorization and uniqueness of prime-power Dirichlet coefficients then exclude every nonzero discrepancy. This gives equality of standard local $L$-factors at every place and entire automorphy of all nontrivial irreducible Artin factors in the stipulated family. Finite explicit-formula certificates also prove real nonvanishing and unconditional automorphy for four explicit fields, of quintic discriminants $8249$, $60061$, $30301$, and $48976$. The tests have nonnegative real part throughout the zero strip; their prime-power terms and integral bounds are certified with directed interval arithmetic. A positive autocorrelation is used for discriminant $30301$, and the wild inertia filtration at two is derived for discriminant $48976$. We also prove that canonical removal of a nonzero positive-definite boundary divisor necessarily loses logarithmic coefficient positivity, at every multiplicity. A finite progression description of the entire divisor would still imply rigidity off the central line; explicit counterexamples delimit what boundary information alone can prove. The general absence of common off-central real zeros remains open here. This manuscript was produced and checked by GPT-6 Astra, with no human review.
\end{abstract}

\section{The published theorem and the refinement}

Throughout, $L/\Q$ is Galois with group $S_5$ and complex conjugation has cycle type $(12)(34)$. No restriction is imposed on ramification or Frobenius at five. Under these assumptions the modularity theorem of \citet[Th\'eor\`eme 0.3]{pillonistroh} supplies the two-dimensional input for Calegari's construction. The replacement is explained in Section~\ref{sec:moderninput}; \citet[Theorem 1.2 and Section 4.2]{kharelarosawiese} make it explicitly for the four-dimensional representation and credit \citet{dwyer2014} with the earlier observation that the local condition at five is superfluous.

\citet[Theorems 1.1--1.2]{calegari} originally also assumed that $L/\Q$ was unramified at five with Frobenius of type $(12)(34)$. His Theorem 1.1 proves automorphy for irreducible representations of dimensions four and six. In dimension five it constructs a tempered cuspidal representation whose local parameters agree with the Artin parameters on a density-one set of primes. The potentially exceptional primes need not initially form a finite set. The published modularity upgrade enlarges the arithmetic family; the central-zero argument below is a separate refinement of the analytic hypothesis.

Let $E=L^{A_5}$, let $F=L^{F_{20}}$ be a sextic resolvent field, where $F_{20}$ is the normalizer of a Sylow $5$-subgroup, and put $H=EF$. Calegari's Theorem 1.2 assumes
\begin{equation}\label{eq:original}
\zeta_H(\sigma)\ne0\qquad(0<\sigma<1,\quad \sigma\in\R).
\end{equation}
It then identifies the automorphic and Artin $L$-functions. The text immediately following that theorem observes that nonvanishing of the particular five-dimensional Artin factor throughout $(0,1)$ would already suffice. We permit a central zero of arbitrary order and replace separate nonvanishing by absence of common real zeros of two entire Hecke functions.

Write $\eta$ for the quadratic character of $E/\Q$. We distinguish the two irreducible five-dimensional representations by choosing $\rh$ with character values
\begin{equation}\label{eq:rhochar}
\chi_{\rh}=(5,1,1,-1,-1,0,1)
\end{equation}
on the conjugacy classes of cycle types
$1,2,2^2,3,4,5,3\cdot2$, in that order. The other representation is $\rh\otimes\eta$. Let $\rho_4$ be the standard four-dimensional representation and $\rho_6=\wedge^2\rho_4$. Let $U=S_2\times S_3$ be the stabilizer of an unordered pair and $U'=S_2\times A_3$ its index-two subgroup. Define
\[
J=L^U,\qquad J'=L^{U'},\qquad [J:\Q]=10,\quad[J':\Q]=20.
\]
The extension $J'/J$ is quadratic.

\begin{theorem}[Common-zero refinement]\label{thm:main}
Let $L/\Q$ be Galois with group $S_5$ and complex conjugation of type $(12)(34)$, and let $\pi$ be the specified tempered cuspidal comparison representation for $\rh$ constructed from Calegari's method using Pilloni--Stroh's modularity theorem. Set
\[
B(s)=\frac{\zeta_H(s)}{\zeta_F(s)},\qquad
C(s)=\frac{\zeta_{J'}(s)}{\zeta_J(s)},
\]
understood as their entire quadratic Hecke continuations. Let $M$ be the completed automorphic-to-Artin comparison quotient defined in \eqref{eq:M}. Unconditionally, if $M$ has a pole of order $r>0$ at any complex point $z$ with $0<\Rea z<1$, then
\begin{equation}\label{eq:mainpolebound}
r\le \min\{\ord_z B,\ord_z C\}.
\end{equation}
This assertion allows arbitrary pole multiplicities and does not require the following nonvanishing hypothesis.

Assume in addition
\begin{equation}\label{eq:RBC}
\bigl(B(\sigma),C(\sigma)\bigr)\ne(0,0)
\qquad\bigl(0<\sigma<1,\ \sigma\ne\tfrac12,\ \sigma\in\R\bigr).
\end{equation}
Then the standard local $L$-factors of $\pi$ and $\rh$ agree at every finite place and at infinity. In particular,
\[
L(s,\pi)=L(s,\rh),\qquad
L(s,\pi\otimes\eta)=L(s,\rh\otimes\eta),
\]
and the two Artin functions are entire. Every nontrivial irreducible complex representation of $\Gal(L/\Q)$ has an entire automorphic $L$-function. Moreover,
\begin{equation}\label{eq:orders}
\ord_{s=1/2}B(s)
=\ord_{s=1/2}L(s,\eta)
+\ord_{s=1/2}L(s,\rh).
\end{equation}
No restriction on the multiplicities of allowed common central zeros is required.
\end{theorem}

\begin{corollary}[Separate nonvanishing]\label{cor:separate}
Under the arithmetic assumptions of Theorem~\ref{thm:main}, all its conclusions hold if either $B$ alone or $C$ alone has no off-central real zero in $(0,1)$. In particular the sufficient condition
\begin{equation}\label{eq:RB}
B(\sigma)\ne0\qquad
\bigl(0<\sigma<1,\ \sigma\ne\tfrac12,\ \sigma\in\R\bigr)
\end{equation}
allows a central zero of arbitrary order.
\end{corollary}

Either separate nonvanishing condition implies \eqref{eq:RBC}.

\begin{remark}[Precise meaning of the local conclusion]\label{rem:local}
The all-place assertion in Theorem~\ref{thm:main} is equality of the standard local factors, including the gamma factors, as in \citet[Lemma 4.10]{calegari}. Equality of a standard $L$-factor alone does not determine a general ramified Weil--Deligne representation. We do not strengthen the conclusion to full ramified parameter equality by that inference.
The weak-to-strong comparison invoked there is formulated in terms of local $L$- and epsilon-factors in \citet[Proposition 4.1]{ramakrishnan}.
\end{remark}

\begin{remark}[An Artin-factor variant]\label{rem:weaker}
The proof also works if the meromorphic Artin function $L(s,\rh)$ and $C(s)$ have no common off-central real zero, where a zero means positive meromorphic order. In particular nonvanishing of $L(s,\rh)$ alone throughout that interval is sufficient. The intrinsically entire formulation \eqref{eq:RBC} avoids assuming Artin holomorphy in advance.
\end{remark}

The difference from \eqref{eq:original} is exact. Zeros belonging to the extraneous factor $\zeta_F$ no longer matter, and central vanishing of $B$ is allowed. The first point follows directly from Calegari's own factorization and observation. The common-zero condition follows from the exterior-square comparison, and the arbitrary central multiplicity is handled by the central-divisor argument. The common-zero hypothesis is logically weaker than either separate nonvanishing hypothesis; no example with an off-central zero is asserted. The hypothesis on complex conjugation remains; the former local condition at five is removed by the published modularity input. Neither condition \eqref{eq:RBC} nor the conclusion is a claim of GRH. Section~\ref{sec:certificate} proves the stronger absence of every real zero for four explicit fields of quintic discriminants $8249$, $60061$, $30301$, and $48976$, giving unconditional examples in this enlarged family.

\subsection*{Boundary divisors beyond the central line}

A further result identifies exactly why canonically removing an outer divisor cannot simply restart the positivity argument. Write \(D_{\rm ar}\) for the arithmetic logarithmic derivative \eqref{eq:Dar}, whose prime-power coefficients satisfy \eqref{eq:bound}.

\begin{theorem}[Canonical removal and complete-divisor rigidity]\label{thm:boundary-method}
\begin{enumerate}[label=(\roman*)]
\item Let \(\alpha>0\), and let \(r\ne0\) be an integer-valued positive-definite function on \(\R_d\) with locally finite ordinary support. Choose its finite coset decomposition
\[
r(t)=\sum_jc_j\one_{\beta_j+h_j\Z}(t)
+\sum_u b_u\one_{\{t_u\}}(t),
\qquad \lambda_j=2\pi/h_j,
\]
and form the canonical factor
\[
P_\alpha(s)=
\prod_j(1-e^{-\lambda_j(s-\alpha-i\beta_j)})^{-c_j}
\prod_u(s-\alpha-it_u)^{-b_u}.
\]
Then \(-P_\alpha'/P_\alpha\) is the Laplace transform of a nonzero nonnegative measure on positive frequencies, but
\[
D_{\rm ar}(s)+P_\alpha'(s)/P_\alpha(s)
\]
cannot be an absolutely convergent generalized Dirichlet series with nonnegative coefficients and positive frequencies. The conclusion allows arbitrary positive \(r(0)\).
\item If the entire signed divisor of the comparison quotient \(M\) in \eqref{eq:M} is a finite integral sum of vertical arithmetic progressions and individual points, all with real parts strictly between zero and one, then \(M\) is constant. The progressions may lie on different vertical lines.
\end{enumerate}
\end{theorem}

The complete proofs are Propositions~\ref{prop:ob-peeling} and \ref{prop:ob-global} in Appendix~\ref{sec:offcentral-obstructions}.

Part (i) concerns canonical elementary factors with their logarithmic constant normalized to zero; it is not an obstruction to every possible reduction. Part (ii) requires the entire divisor, whereas the Abel limit gives only the outer boundary. The counterexamples in Appendix~\ref{sec:offcentral-obstructions} separately preserve exact prime Euler factors or global reflection and finite order, and state precisely which remaining hypothesis each omits. These method results do not reduce the remaining interval for an off-central comparison pole.

\section*{Proof overview}

Calegari's unresolved primes contribute a positive Euler quotient
$\prod_{p\in\Omega}((1+p^{-s})/(1-p^{-s}))^2$.
Taking exterior squares preserves the same exceptional quotient, and $\wedge^2\rh=\rho_4\otimes\eta\oplus\rho_6$ is the Artin representation of $C$. The partial exterior-square function of the cuspidal rank-five comparison is entire. Every pole of the original quotient is therefore a zero of $C$, as well as of $B$, with at least its pole order. If $\Omega$ were infinite, Landau's theorem would give a real pole. Condition \eqref{eq:RBC} puts that pole at $1/2$; the functional equation then confines the entire divisor of the completed quotient to the critical line. The residues of its logarithmic derivative are pointwise limits of positive sums of characters. They form a bounded integer-valued positive-definite function on $\R$ equipped with the discrete topology. Cohen's theorem forces its level sets into a coset ring. Ordinary local finiteness of the divisor then reduces those cosets to arithmetic progressions and isolated points.

The remaining argument uses analytic continuation in a concrete way. A finite progression divisor is realized by a finite product of exponential factors; an isolated divisor is realized by a rational function. Hadamard factorization shows that these account for the quotient up to an exponential. Gamma recurrence eliminates the isolated part. Finally, the progression factors would produce coefficients with exponential growth on finitely many frequency lattices. Actual local Artin and tempered automorphic parameters give bounded prime-power coefficients. Unique factorization of rational integers makes these two descriptions compatible only when the discrepancy vanishes.

The explicit examples use separate positive tests in the classical explicit formula. Dividing a triangular function by $\cosh(x/2)$ makes the real part of its Fourier transform nonnegative throughout the full critical strip. Any real zero would contribute a fixed positive minimum to the zero sum, including at the center because the Hecke functional equation has sign $+1$. Exact prime-power traces and interval bounds show that the whole zero sum is smaller than that minimum.

For the boundary-removal result, Bochner's measure splits into mutually singular Haar components indexed by rationally commensurate height lattices, together with the Haar component from isolated points. Positivity makes every component nonnegative. The canonical logarithmic derivative therefore has a continuous positive part or positive lattice coefficients with exponential growth. Subtracting either part from bounded prime-power coefficients loses the positivity required by the next Abel argument.

\section{Fields, characters, and the Hecke quotients}

The subgroup $F_{20}$ is the affine linear group of the line over $\mathbb F_5$. Its twenty elements consist of the identity, four $5$-cycles, five double transpositions, and ten $4$-cycles. Thus $F_{20}$ contains odd permutations and
\[
F_{20}\cap A_5=D_{10},\qquad [H:F]=2.
\]
Here $D_{10}$ has order ten. Frobenius reciprocity and the character table give
\begin{equation}\label{eq:induction}
\Ind_{F_{20}}^{S_5}\one=\one\oplus(\rh\otimes\eta),
\qquad
\Ind_{D_{10}}^{S_5}\one
=\one\oplus\eta\oplus\rh\oplus(\rh\otimes\eta).
\end{equation}
For example, averaging \eqref{eq:rhochar} over $F_{20}$ gives
$(5+5-10)/20=0$, while averaging its sign twist gives
$(5+5+10)/20=1$. Together with dimensions and the other character averages, these are exactly the decompositions in \eqref{eq:induction}.

Artin formalism, including ramified Euler factors, yields
\begin{equation}\label{eq:factorization}
\zeta_F(s)=\zeta_{\Q}(s)L(s,\rh\otimes\eta),\qquad
\zeta_H(s)=\zeta_F(s)L(s,\eta)L(s,\rh).
\end{equation}
Equivalently,
\begin{equation}\label{eq:B}
B(s)=L_F(s,\chi_{H/F})=L(s,\eta)L(s,\rh),
\end{equation}
where $\chi_{H/F}$ is the nontrivial quadratic Hecke character of $F$. The first equality proves that $B$ is entire. The second is initially an equality of meromorphic functions, and therefore gives equality of signed orders at every point. Formula \eqref{eq:orders} consequently needs no new multiplicity calculation after holomorphy has been obtained.

Since complex conjugation lies in $A_5$, the field $E$ is real quadratic. Its two real places lie above the real place of $\Q$. The Artin parameter of $\rh$ at infinity has three eigenvalues $+1$ and two eigenvalues $-1$: its dimension is five and its trace on a double transposition is one. Thus
\begin{equation}\label{eq:gamma}
L_\infty(s,\rh)=\Gr(s)^3\Gr(s+1)^2
=\Gr(s)\Gc(s)^2,
\end{equation}
with $\Gr(s)=\pi^{-s/2}\Gamma(s/2)$ and
$\Gc(s)=2(2\pi)^{-s}\Gamma(s)$.

\begin{proposition}[The second quadratic Hecke function]\label{prop:secondHecke}
There are identities of representations and of all-place Artin functions
\begin{align}
\wedge^2\rh
&=(\rho_4\otimes\eta)\oplus\rho_6
=\Ind_{S_2\times S_3}^{S_5}\operatorname{sgn}_{S_3},\label{eq:exteriorHecke}\\
C(s)&=L_J(s,\chi_{J'/J})
=L(s,\rho_4\otimes\eta)L(s,\rho_6).\label{eq:CHecke}
\end{align}
In particular $C$ is entire, and
\[
C_\infty(s)=\Gr(s)^4\Gr(s+1)^6.
\]
\end{proposition}

\begin{proof}
The character formula
\[
\chi_{\wedge^2\rh}(g)
=\tfrac12\bigl(\chi_{\rh}(g)^2-\chi_{\rh}(g^2)\bigr)
\]
gives $(10,-2,-2,1,0,0,1)$ on the ordered classes in \eqref{eq:rhochar}. This equals the character of $\rho_4\otimes\eta\oplus\rho_6$, since the respective characters of $\rho_4$ and $\rho_6$ are
\[
(4,2,0,1,0,-1,-1),\qquad(6,0,-2,0,0,1,0).
\]

Let $P=\one\oplus\rho_4$ be the permutation representation on five letters. Its exterior-square basis is indexed by unordered pairs, and the pair stabilizer acts on a basis vector by $\operatorname{sgn}_{S_2}$. Hence
\[
\wedge^2P=\rho_4\oplus\rho_6
=\Ind_{S_2\times S_3}^{S_5}\operatorname{sgn}_{S_2}.
\]
Twisting by $\eta$ cancels the sign on the chosen pair and leaves the complementary triple sign. Also $\rho_6\otimes\eta=\rho_6$, as its character vanishes on odd classes. This proves \eqref{eq:exteriorHecke}. The kernel of $\operatorname{sgn}_{S_3}$ is $U'$. Induction in stages and Artin formalism therefore identify $\zeta_{J'}/\zeta_J$ with the nontrivial quadratic Hecke function in \eqref{eq:CHecke}, including ramified factors. Hecke's theorem makes it entire. Its parameter at complex conjugation has dimension ten and trace $-2$, giving four positive and six negative eigenvalues and the displayed gamma factor.
\end{proof}

Both completed Hecke functions $B$ and $C$ have functional equation with sign $+1$, since each is a quotient of completed Dedekind zeta functions for a quadratic extension. Their real zeros in $(0,1)$ are therefore reflected about $1/2$ with the same orders. Equivalently, the common-zero condition \eqref{eq:RBC} can be imposed only on $(1/2,1)$.

\section{The exact Calegari comparison input}

\subsection{The modularity input without a condition at five}\label{sec:moderninput}

\begin{proposition}[Published scope upgrade]\label{prop:scopeupgrade}
For every $S_5$-extension with complex conjugation of type $(12)(34)$, Calegari's two odd binary icosahedral representations over $E$ are modular. Consequently his four- and six-dimensional automorphy results and his five-dimensional comparison construction apply without a condition at five.
\end{proposition}

\begin{proof}
Since complex conjugation lies in $A_5$, the quadratic field $E$ is real. Calegari's Lemma 3.1 constructs a finite cyclic central extension of the $S_5$ quotient by lifting its projective realization in $\operatorname{PGL}_2(\mathbb F_5)$, using Tate's lifting theorem. The finite-group construction in his Section 2 then gives two faithful complex representations
\[
\varrho,\ \varrho^\iota:G_E\longrightarrow\GL_2(\C)
\]
with projective image $A_5$. They have finite image and are irreducible: a reducible finite-image two-dimensional representation is a sum of characters and cannot have this projective image.

At either real place of $E$, the image of complex conjugation is a non-scalar involution, since its projective image is the nonidentity double transposition. Its eigenvalues are $1,-1$, so its determinant is $-1$. Both representations are therefore totally odd. \citet[Th\'eor\`eme 0.3]{pillonistroh} attaches a cuspidal Hilbert eigenform of parallel weight one to each. This theorem has no residual, ordinarity, conductor, or auxiliary-prime hypothesis. It replaces the appeal to Sasaki and the condition on $\operatorname{Frob}_5$ in \citet[Lemma 3.2]{calegari}.

For completeness, every subsequent dependency can be traced explicitly. Calegari's Section 3 obtains the four-dimensional representation from the Asai transfer and the six-dimensional representation from its exterior square, or from quadratic automorphic induction of the adjoint of $\varrho$. His Section 4 uses the symmetric fourth power and quadratic descent, then the descended symmetric cube, its symplectic structure, and its exterior square. Lemmas 4.5--4.6 use Rankin--Selberg poles and Chebotarev to determine the similitude character. The local tables and Theorem 4.7 are finite representation calculations. None of these steps uses the splitting or ramification of five after the two-dimensional modularity input has been established.

All ramified places, including five if necessary, are included in a finite set. Base change gives finite-image local parameters, so the bad local eigenvalues satisfy the same bounds as in Calegari's argument. The two real places of $E$ give the archimedean comparison; the sign character is even. Thus Lemmas 4.8--4.10, including the finite-place twisting argument and the gamma-factor comparison, retain their hypotheses. This proves the claimed extension of the entire comparison construction.
\end{proof}

The infinity condition is still needed for this route. Trivial complex conjugation gives an even projective lift, while a transposition makes $E$ imaginary; neither case falls under the stated totally real, totally odd modularity theorem. The determinant calculation above also fixes the sign in the corresponding argument in \citet[Section 4.2]{kharelarosawiese}.

\subsection{The precise comparison package}

We isolate the established automorphic input to avoid replacing a density-one assertion by an all-but-finite assertion. The latter would assume the point to be proved.

\begin{proposition}[Calegari's comparison package]\label{prop:package}
There are a finite set $S$ of rational primes and a set $\Omega$ of primes disjoint from $S$ such that:
\begin{enumerate}[label=(\alph*)]
\item $\pi$ is tempered and cuspidal on $\GL_5/\Q$; outside $S\cup\Omega$ its local parameters equal those of $\rh$;
\item at $p\in\Omega$, the automorphic parameters are
$\{1,1,1,i,-i\}$, whereas the Artin parameters are
$\{1,-1,-1,i,-i\}$;
\item $\sum_{p\in\Omega}p^{-1}<\infty$;
\item if $\Omega$ is finite, the standard local factors agree at every place, and the same conclusion holds after twisting by $\eta$.
\end{enumerate}
\end{proposition}

The sources are \citet[Theorem 1.1, Theorem 4.7, Lemmas 4.8--4.10]{calegari}. Part (c) is the monotone-convergence consequence of the bounded prime sum in the proof of Theorem 4.7; it is stronger than its stated density-zero conclusion. Lemma 4.9 supplies the Landau argument, and Lemma 4.10 removes a finite exceptional set.
More explicitly, the logarithm of the Rankin--Selberg comparison quotient used there is bounded as real $s\downarrow1$, since its numerator and denominator both have simple poles. Its prime-term calculation gives $8\sum_{p\in\Omega}p^{-s}+O(1)$, with the higher prime powers uniformly summable. Hence these positive prime sums are bounded, and monotone convergence gives (c). Only this boundedness is needed here; the density statement alone would not imply (c).

For clarity fix the reciprocal-polynomial convention
\[
L_p(s,\tau)=P_{\tau,p}(p^{-s})^{-1}.
\]
The two polynomials at an exceptional prime are
\[
P_{\pi,p}(T)=(1+T^2)(1-T)^3,
\qquad P_{\rh,p}(T)=(1+T^2)(1-T)(1+T)^2.
\]
Thus the quotient is $P_{\rh,p}/P_{\pi,p}$, and
\begin{equation}\label{eq:Q}
Q_\Omega(s)=\prod_{p\in\Omega}
\left(\frac{1+p^{-s}}{1-p^{-s}}\right)^2.
\end{equation}
The quotient in \eqref{eq:Q} is formed from the inverse local polynomials in this fixed convention.

Let $\Lambda(s,\tau)=N_\tau^{s/2}L_\infty(s,\tau)L(s,\tau)$ denote the conductor-completed function, and put
\begin{equation}\label{eq:M}
M(s)=\frac{\Lambda(s,\pi)}{\Lambda(s,\rh)}
=A^{s/2}G(s)Q_S(s)Q_\Omega(s),
\qquad A=\frac{N_\pi}{N_{\rh}}>0.
\end{equation}
The finite quotient $Q_S$ is a product over $p\in S$. Calegari's construction gives gamma factors made from $\Gr(s)$ and $\Gc(s)$, of equal total degree five, so for some integer $d$,
\begin{equation}\label{eq:G}
G(s)=\left(\frac{\Gr(s)}{\Gr(s+1)}\right)^d.
\end{equation}
In fact $d=0$ follows directly from real-place base change; we permit $d$ at this stage because the proof below eliminates it without that additional identification.

All local inverse roots at finite places have modulus at most one. The finite quotient $Q_S$, and the gamma quotient $G$, are therefore holomorphic and nonzero for $\Rea s>0$. In the particular construction Calegari obtains roots of unity. Only the modulus bound is needed in the coefficient argument below.

The automorphic representation $\pi$ is self-dual. Indeed every unramified parameter multiset is invariant under inversion, both at the matching primes and at the exceptional primes. Strong multiplicity one, applied to the two already automorphic representations $\pi$ and $\pi^\vee$, gives $\pi\simeq\pi^\vee$. Consequently,
\begin{equation}\label{eq:FE}
M(s)=\epsilon M(1-s)
\end{equation}
for a nonzero constant $\epsilon$. The exact value of $\epsilon$ is unnecessary. Brauer induction expresses the completed Artin function as a finite product and quotient of completed Hecke functions. Together with the standard analytic properties of the cuspidal numerator, it follows that $M$ is meromorphic of order at most one. Here and below this may be understood in the Nevanlinna sense, or by writing the functions as quotients of entire functions of order at most one.

Finally define
\begin{align}
D_\Omega(s)&=-\frac{Q_\Omega'}{Q_\Omega}(s)
=4\sum_{p\in\Omega}\sum_{\substack{k\ge1\\k\ \mathrm{odd}}}
(\log p)p^{-ks},\label{eq:Dpositive}\\
D_{\mathrm{ar}}(s)&=-\frac{(Q_SQ_\Omega)'}{Q_SQ_\Omega}(s)
=\sum_p\sum_{k\ge1}a(p,k)(\log p)p^{-ks}.
\label{eq:Dar}
\end{align}
Each $a(p,k)$ is a difference of traces of powers of at most five local inverse roots. Therefore
\begin{equation}\label{eq:bound}
|a(p,k)|\le10.
\end{equation}
For $p\in\Omega$, these coefficients equal $4$ for odd $k$ and $0$ for even $k$. No bad-prime term can alter them because $S\cap\Omega=\varnothing$.

\subsection{Exterior-square control of every comparison pole}

\begin{proposition}[Common-zero divisor bound]\label{prop:commonpole}
If $z$ lies in $0<\Rea z<1$ and $M$ has a pole there of order $r>0$, then
\begin{equation}\label{eq:commonorder}
r\le\min\{\ord_zB,\ord_zC\}.
\end{equation}
This holds at arbitrary multiplicity and at nonreal as well as real points.
\end{proposition}

\begin{proof}
Enlarge the finite set $S$ to contain every ramified place of the functions under consideration, and omit infinity as well. Write
\[
\mathcal E^S(s)=L^S(s,\pi,\wedge^2),\qquad
C^S(s)=\prod_{p\notin S}L_p(s,C).
\]
The representation $\pi$ is unitary cuspidal of odd rank five. By \citet[Theorem 6.2]{belt}, the partial exterior-square function $\mathcal E^S$ is entire. The twisting character in that theorem is trivial, and $S$ contains all its required places. Completed odd-rank exterior-square entirety is also proved by \citet[Theorem 3.5]{kim1999}. Only the partial holomorphy statement is needed here.

At a good matching prime the exterior-square local factors agree. At $p\in\Omega$, taking the ten pairwise products of the two spectra in Proposition~\ref{prop:package} gives multiplicities
\[
\begin{array}{c|rrrr}
\text{eigenvalue}&1&i&-1&-i\\\hline
\wedge^2\rho_{5,p}&2&3&2&3\\
\wedge^2\pi_p&4&3&0&3.
\end{array}
\]
Thus, with $T=p^{-s}$,
\[
\frac{L_p(s,\pi,\wedge^2)}{L_p(s,C)}
=\frac{(1-T)^2(1+T)^2(1+T^2)^3}
{(1-T)^4(1+T^2)^3}
=\left(\frac{1+T}{1-T}\right)^2.
\]
The reciprocal-polynomial orientation is the same as in the original comparison. Consequently
\begin{equation}\label{eq:exteriorpartial}
\mathcal E^S(s)=C^S(s)Q_\Omega(s)
\end{equation}
first for $\Rea s>1$, and then meromorphically everywhere.

Every omitted finite factor of $C$ is the Artin factor of a finite-image representation. Its inverse roots on inertia invariants have modulus one. Hence the finite polynomial product $C^S/C$ is holomorphic and nowhere zero for $\Rea s>0$, and $\ord_zC^S=\ord_zC$ in that half-plane. Equation \eqref{eq:M} likewise shows that $\ord_zQ_\Omega=\ord_zM$ there: the conductor exponential, $G$, and $Q_S$ have no divisor. This accounts for all finite and archimedean corrections without identifying any ramified exterior-square factor of $\pi$.

At the stipulated pole, taking orders in \eqref{eq:exteriorpartial} yields
\[
0\le\ord_z\mathcal E^S=\ord_zC-r.
\]
For the other inequality, the standard completed quotient and the absence of gamma or conductor divisors give
\[
\ord_zL(s,\rh)=\ord_zL(s,\pi)+r.
\]
Using \eqref{eq:B} and the entirety of the cuspidal standard function and the nontrivial quadratic Dirichlet function,
\[
\ord_zB=\ord_zL(s,\eta)+\ord_zL(s,\pi)+r\ge r.
\]
No holomorphy of $L(s,\rh)$ was assumed in this order calculation.
\end{proof}

\begin{corollary}[Orders in the pair-stabilizer fields]\label{cor:pairfieldorders}
At a pole as in Proposition~\ref{prop:commonpole},
\[
\ord_z\zeta_J\ge r,\qquad \ord_z\zeta_{J'}\ge2r.
\]
\end{corollary}

\begin{proof}
The pair permutation character is $\one\oplus\rho_4\oplus\rh$, so
\[
\zeta_J(s)=\zeta_{\Q}(s)L(s,\rho_4)L(s,\rh).
\]
The first two factors are holomorphic in the open strip, and the proof of Proposition~\ref{prop:commonpole} gives $\ord_zL(s,\rh)\ge r$. Finally $\zeta_{J'}=\zeta_J C$, and both orders on the right are at least $r$. These are consequences of the same comparison-pole bound, not a separate zero-exclusion theorem.
\end{proof}

\section{Landau's theorem and the critical-line divisor}

\begin{lemma}\label{lem:abscissa}
If $\Omega$ is infinite, the positive Euler product $Q_\Omega$ has abscissa of convergence $\alpha\in(0,1)$ and a pole at $s=\alpha$. Under \eqref{eq:RBC}, $\alpha=1/2$.
\end{lemma}

\begin{proof}
The expansion
\[
\frac{(1+x)^2}{(1-x)^2}=1+4\sum_{n\ge1}nx^n
\]
has nonnegative coefficients. Thus $Q_\Omega$ is an ordinary Dirichlet series with nonnegative coefficients. By Proposition~\ref{prop:package}(c) its Euler product and Dirichlet series converge at $s=1$. Its meromorphic continuation follows from \eqref{eq:M}. Landau's theorem implies that a finite abscissa of convergence is a real singularity and hence a pole; see \citet[Lemma 4.9]{calegari}. Convergence at $1$ rules out a pole there, so $\alpha<1$.

The abscissa cannot be negative, since even one nonempty Euler factor has a divergent nonnegative Dirichlet expansion at $0$. If $\alpha=0$, then for each finite $T\subset\Omega$ and real $s>0$,
\[
Q_\Omega(s)\ge\prod_{p\in T}
\left(\frac{1+p^{-s}}{1-p^{-s}}\right)^2
\asymp_T s^{-2|T|}\qquad(s\downarrow0).
\]
Meromorphicity would force a pole at zero of order at least $2|T|$ for every finite $T$, which is impossible. Hence $\alpha>0$.

At a real pole of $Q_\Omega$ in $(0,1)$, the factors $A^{s/2}GQ_S$ are nonzero and regular, so $M$ has a pole. Proposition~\ref{prop:commonpole} makes it a common zero of $B$ and $C$, with at least the same order in each. Condition \eqref{eq:RBC} therefore gives $\alpha=1/2$.
\end{proof}

For $\Rea s>\alpha$, absolute convergence of the product also gives absolute convergence of its Euler logarithm and of \eqref{eq:Dpositive}. This follows directly by comparing to $\sum_{p\in\Omega}p^{-\sigma}$ and moving to a slightly smaller real part to absorb $\log p$. Consequently $M$ is holomorphic and nonzero for $\Rea s>1/2$. The functional equation \eqref{eq:FE} reflects this statement to $\Rea s<1/2$. We have proved:

\begin{corollary}\label{cor:line}
If $\Omega$ is infinite and \eqref{eq:RBC} holds, every zero and pole of $M$ lies on $\Rea s=1/2$. Its pole order at $1/2$ is a finite positive integer.
\end{corollary}

\section{Integer-valued positive definiteness and Cohen's theorem}

Write $\R_d$ for the additive group of real numbers with the discrete topology. The support of a function on this group is nevertheless called locally finite below only with respect to the usual topology of $\R$.

\begin{lemma}[Residue positivity]\label{lem:PD}
Under the hypotheses of Corollary~\ref{cor:line}, define
\[
r(t)=-\ord_{s=1/2+it}M(s),\qquad t\in\R.
\]
Then $r$ is an integer-valued positive-definite function on $\R_d$, its support is locally finite, and
\[
|r(t)|\le r(0)=m<\infty.
\]
The sign convention assigns positive values to poles.
\end{lemma}

\begin{proof}
The gamma and finite factors in \eqref{eq:M} have no divisor on $\Rea s=1/2$. For every real $t$, the elementary residue formula for the logarithmic derivative gives
\begin{equation}\label{eq:Abel}
r(t)=\lim_{\delta\downarrow0}\delta
D_\Omega(1/2+\delta+it).
\end{equation}
This includes regular points, where the limit is zero. For each $\delta>0$, the function before the limit is an absolutely convergent positive sum of characters $t\mapsto e^{-itk\log p}$. In particular, for arbitrary complex $z_1,\dots,z_q$ and real $t_1,\dots,t_q$,
\[
\sum_{i,j}z_i\overline{z_j}\,
\delta D_\Omega(1/2+\delta+i(t_i-t_j))\ge0.
\]
Taking the limit proves positive definiteness of $r$. Two-by-two principal matrices give $|r(t)|\le r(0)$. Meromorphic orders are integers, and the divisor of a nonzero meromorphic function is locally finite. The value $m=r(0)>0$ is the central pole order from Corollary~\ref{cor:line}.
\end{proof}

Ordinary continuity is not asserted in Lemma~\ref{lem:PD}. For instance, a function supported on a discrete set of real heights is typically discontinuous at each nonzero atom. The topology $\R_d$ is what permits Bochner's theorem to be used here.

\begin{lemma}[The coset consequence]\label{lem:cosets}
Let $r:\R_d\to\Z$ be positive definite and have locally finite support in the ordinary real topology. Then there are finite index sets, integers $c_j,b_t$, real $\beta_j$, and positive real $\alpha_j$ such that
\begin{equation}\label{eq:comb}
r(x)=\sum_{j=1}^Jc_j\one_{\beta_j+\alpha_j\Z}(x)
+\sum_{t\in T}b_t\one_{\{t\}}(x).
\end{equation}
No uniform separation assumption is required.
\end{lemma}

\begin{proof}
Bochner's theorem for the locally compact discrete abelian group $\R_d$ represents $r$ as the Fourier--Stieltjes transform of a finite positive measure on its compact Pontryagin dual. Its total mass is $r(0)$, so $r$ belongs to the Fourier--Stieltjes algebra $B(\R_d)$. The bound $|r(x)|\le r(0)$ and integrality make its range finite.

For each nonzero value $v$ in this range, polynomial interpolation on the range gives a polynomial $P_v$ for which $P_v(r)=\one_{\{r=v\}}$. The algebra $B(\R_d)$ is closed under products and contains constants. Hence the indicator of this level set is a Fourier--Stieltjes transform. Its values are zero and one, so its representing measure is idempotent under convolution. Cohen's idempotent theorem places the level set in the coset ring of $\R_d$ \citep{cohen}; a primary modern statement is \citet[Theorem 1.1]{greensanders}.

Every nonzero level set is locally finite. By \citet[Theorem 3, Section 3.2]{kolountzakis}, a discrete element of the abstract coset ring belongs to the ring generated by ordinary discrete cosets. In one real dimension those cosets are singletons and sets $\beta+\alpha\Z$. Expanding finite unions, intersections, and differences by inclusion--exclusion gives an integer linear combination of their indicators: an intersection of finitely many cosets, if nonempty, is again a coset, and the intersection of discrete subgroups remains discrete. Summing these expressions with coefficient $v$ proves \eqref{eq:comb}.
\end{proof}

The discrete-coset step is substantive. Abstract subgroups of $\R_d$ need not be cyclic; many are dense in $\R$. One cannot replace them by lattices merely because the ambient group is the real line. The cited theorem uses the ordinary discreteness of the level set to remove the dense cosets, including possible Boolean differences. Only after that reduction does the classification of discrete subgroups of $\R$ apply.

\section{Recovering a meromorphic quotient from its divisor}

Assume for contradiction that $\Omega$ is infinite. Apply Lemma~\ref{lem:cosets} to the function in Lemma~\ref{lem:PD}. Put
\[
\lambda_j=\frac{2\pi}{\alpha_j},\qquad \rho_t=\tfrac12+it.
\]
The meromorphic function
\begin{equation}\label{eq:P}
P(s)=\prod_{j=1}^J
\left(1-e^{-\lambda_j(s-1/2-i\beta_j)}\right)^{-c_j}
\prod_{t\in T}(s-\rho_t)^{-b_t}
\end{equation}
has the same divisor as $M$. Indeed the zeros of the $j$th exponential factor are precisely $1/2+i(\beta_j+\alpha_j\Z)$, and all are simple. Repeated or intersecting progressions cause no difficulty because orders add, exactly as the coefficients in \eqref{eq:comb} do.

Both $M$ and $P$ are meromorphic of order at most one. Therefore $M/P$ is a nonvanishing entire function of order at most one. The zero-free case of Hadamard factorization gives
\begin{equation}\label{eq:Hadamard}
M(s)=e^{as+b}P(s)
\end{equation}
for constants $a,b\in\C$. This use of order is essential; a zero-free entire function with no order restriction need not be an exponential of a linear polynomial.

Taking negative logarithmic derivatives in \eqref{eq:Hadamard}, we obtain
\begin{equation}\label{eq:modelD}
-\frac{M'}M(s)=-a+R(s)+C(s),
\end{equation}
where
\begin{align}
R(s)&=\sum_{t\in T}\frac{b_t}{s-\rho_t},\label{eq:R}\\
C(s)&=\sum_{j=1}^Jc_j\lambda_j
\sum_{n\ge1}e^{n\lambda_j/2}e^{in\lambda_j\beta_j}e^{-n\lambda_js}.
\label{eq:C}
\end{align}
The second series is absolutely convergent in $\Rea s>1/2$. Its sign follows from differentiating $(1-e^{-\lambda(s-\rho)})^{-c}$: its negative logarithmic derivative is $c\lambda/(e^{\lambda(s-\rho)}-1)$.

\begin{lemma}[Eliminating isolated factors and gamma mismatch]\label{lem:gamma}
In \eqref{eq:modelD}, $R=0$, the integer $d$ in \eqref{eq:G} is zero, and
\begin{equation}\label{eq:equalD}
D_{\mathrm{ar}}(s)=C(s)
\end{equation}
in a right half-plane.
\end{lemma}

\begin{proof}
Let $h=-G'/G$. Equation \eqref{eq:M} gives
\[
-M'/M=-\tfrac12\log A+h+D_{\mathrm{ar}}.
\]
As $s\to+\infty$ along the real axis, both $h(s)$ and $R(s)$ tend to zero, while $D_{\mathrm{ar}}(s)$ and $C(s)$ decay exponentially. Comparing constants first gives $a=(\log A)/2$, and then
\begin{equation}\label{eq:RminusH}
R(s)-h(s)=O(e^{-cs})
\end{equation}
for some $c>0$. The Dirichlet series decays exponentially by absolute convergence at any fixed real point greater than one; the same assertion for $C$ follows from its finitely many positive fundamental frequencies.

The recurrence $\Gr(s+2)=s\Gr(s)/(2\pi)$ implies
\[
h(s+2)-h(s)=-\frac d s+\frac d{s+1}.
\]
Subtracting the shift by two of \eqref{eq:RminusH} shows that the rational function
\begin{equation}\label{eq:rationalshift}
R(s+2)-R(s)+\frac d s-\frac d{s+1}
\end{equation}
is exponentially small along the positive real axis. A nonzero rational function has a nonzero leading algebraic asymptotic there. Hence \eqref{eq:rationalshift} vanishes identically.

Every pole of $R(s)$ has real part $1/2$, and every pole of $R(s+2)$ has real part $-3/2$. The two displayed rational gamma terms have poles only at $0$ and $-1$. These locations are disjoint, forcing $d=0$. We then have $R(s+2)=R(s)$. A rational periodic function tending to zero at infinity is identically zero. Substitution in the logarithmic derivative identities proves \eqref{eq:equalD}.
\end{proof}

\section{Prime-power coefficients exclude the progression divisor}

\begin{lemma}[Finite phases cannot decay]\label{lem:phases}
If $u(n)=\sum_{j=1}^qz_je^{in\theta_j}$, where the phases are distinct modulo $2\pi$, and $u(n)\to0$ as $n\to\infty$, then $z_1=\cdots=z_q=0$.
\end{lemma}

\begin{proof}
The Cesaro mean of $|u(n)|^2$ tends to $\sum_j|z_j|^2$, since the averages of every nonconstant phase $e^{in(\theta_j-\theta_k)}$ tend to zero. If $u(n)\to0$, the same mean tends to zero.
\end{proof}

\begin{proposition}[Arithmetic rigidity]\label{prop:rigidity}
The identity \eqref{eq:equalD}, together with \eqref{eq:Dar} and the uniform coefficient bound \eqref{eq:bound}, forces $D_{\mathrm{ar}}=0$.
\end{proposition}

\begin{proof}
Partition the finite set of frequencies $\lambda_j$ into classes under the equivalence relation $\lambda_j/\lambda_k\in\Q$. Within a class choose $\lambda>0$ such that $\lambda_j=m_j\lambda$ with positive integers $m_j$. This is possible by clearing the finitely many rational denominators. Positive integer multiples of frequencies from distinct classes never coincide.

At the frequency $n\lambda$, the coefficient contributed by this class in \eqref{eq:C} is
\begin{equation}\label{eq:un}
e^{n\lambda/2}u(n),\qquad
u(n)=\sum_{j:m_j\mid n}c_jm_j\lambda e^{in\lambda\beta_j}.
\end{equation}
The elementary identity
\[
\one_{m\mid n}=\frac1m\sum_{r=0}^{m-1}e^{2\pi irn/m}
\]
shows that $u(n)$ is a finite sum of unit-modulus phases raised to the $n$th power.

The union of the two sets of frequencies in \eqref{eq:equalD} is locally finite, and both generalized Dirichlet series are absolutely convergent in a common right half-plane. Uniqueness identifies their coefficients at every frequency. For completeness, this uniqueness follows by taking the least frequency with a nonzero difference, multiplying by its exponential, and sending real $s$ to infinity; dominated convergence at a fixed point of absolute convergence handles the tail. Iteration proves that every difference coefficient is zero.

The set $\lambda\N$ can meet the prime-power logarithms $\{k\log p:p\text{ prime},\ k\ge1\}$ for at most one distinct prime $p$. Otherwise $n\lambda=k\log p$ and $n'\lambda=k'\log q$ with $p\ne q$ would imply
$p^{kn'}=q^{k'n}$, contrary to unique factorization. Consequently the coefficient on the left side of \eqref{eq:equalD} at $n\lambda$ is either always zero, or bounded in magnitude by $10\log p$ for a single fixed prime $p$. Equations \eqref{eq:equalD}--\eqref{eq:un} therefore imply
\[
u(n)=O(e^{-n\lambda/2}).
\]
Lemma~\ref{lem:phases} forces this class's coefficients to vanish identically. Every class has the same property, so $C=0$ and hence $D_{\mathrm{ar}}=0$.
\end{proof}

\begin{proof}[Proof of Theorem~\ref{thm:main}]
If $\Omega$ were infinite, Sections 4--7 would apply. Proposition~\ref{prop:rigidity} would give $D_{\mathrm{ar}}=0$. But the coefficient of $p^{-s}$ in \eqref{eq:Dar} is $4\log p$ for every $p\in\Omega$, an immediate contradiction. Thus $\Omega$ is finite.

Proposition~\ref{prop:package}(d), which is \citet[Lemma 4.10 and the final proof of Theorem 1.2]{calegari}, now gives equality of standard local factors at every place, including those in the finite exceptional set and the archimedean place. It gives the same conclusion after the quadratic twist. Standard $L$-functions of these cuspidal representations are entire. Calegari's Theorem 1.1 covers the two four-dimensional representations and the six-dimensional representation; Hecke theory covers the nontrivial one-dimensional character. These are all nontrivial irreducible representations of $S_5$. Finally \eqref{eq:B} gives \eqref{eq:orders} with nonnegative orders. The proof used the finiteness of $m=r(0)$, but no upper bound on it.
\end{proof}

\section{What Booker's degree-zero theorems contribute}

The comparison can also be expressed in the language of $L$-data developed by \citet{booker2015}. It is important to retain his normalization and his meaning of positivity.

An $L$-datum is a triple $\mathcal F=(f,K,m)$ satisfying growth, regularity, and explicit-formula axioms. Its function is
\begin{equation}\label{eq:bookerL}
L_{\mathcal F}(s)
=\exp\left(\sum_{n\ge2}\frac{f(n)}{\log n}n^{1/2-s}\right)
=\sum_{n\ge1}a_{\mathcal F}(n)n^{-s}.
\end{equation}
The degree is $2\lim_{x\downarrow0}xK(x)$. The datum is \emph{positive} if $m(z)<0$ at only finitely many points. For completed meromorphic functions, this says that there are only finitely many poles. It is not the assertion that the coefficients $f(n)$ are nonnegative.

\citet[Theorem 1.6]{booker2015} proves the following equivalent criteria for a datum to vanish:
\begin{equation}\label{eq:bookercriteria}
\begin{gathered}
\mathcal F=0,\qquad
\sum_{n\ge2}\frac{|f(n)|}{\log n}<\infty,\qquad
\sum_{n\ge1}\frac{|a_{\mathcal F}(n)|}{\sqrt n}<\infty,\\
L_{\mathcal F}\text{ is a ratio of Dirichlet polynomials},\qquad
\sum_{\substack{z\in\supp m\\|\Rea z|\le T}}|m(z)|=o(T).
\end{gathered}
\end{equation}
His Theorem 1.7 classifies positive data of degree less than $5/3$; the degree-zero case says that a positive datum of degree zero is zero. These statements apply to Artin data without first assuming Artin holomorphy; see his Examples 1.4(2)--(3).

For the present quotient the difference datum has degree zero. Its coefficients at good primes are
\[
f_R(p^k)=
\begin{cases}
4(\log p)p^{-k/2},&p\in\Omega,\ k\text{ odd},\\
0,&\text{otherwise}.
\end{cases}
\]
The remaining coefficients lie on powers of the finitely many bad primes. If a nonzero quotient with the central pole existed, Booker's theorem would imply that its completed divisor has infinitely many poles and infinitely many zeros: otherwise the difference datum, or its negative, would be positive of degree zero and hence zero. The same theorem rules out an $o(T)$ total variation of that signed divisor. These are obstructions to a nonzero discrepancy, not a proof that the discrepancy is zero.

\begin{proposition}[Euler positivity is different from datum positivity]\label{prop:counterexample}
Let $\chi_5$ be the even primitive quadratic character modulo five. The degree-zero datum associated with $\zeta(s)/L(s,\chi_5)$ has $K=0$ and
\[
f(n)=(1-\chi_5(n))\frac{\Lambda(n)}{\sqrt n}\ge0
\qquad(n>1),
\]
but it is not positive in Booker's sense.
\end{proposition}

\begin{proof}
The gamma kernels cancel because both degree-one factors are even. The displayed coefficients follow by subtracting logarithmic derivatives. At $n=2$ the coefficient is nonzero since $\chi_5(2)=-1$. If the datum were positive, Booker's degree-zero theorem would make it zero, a contradiction. Applying the same argument to its negative shows that the completed quotient has infinitely many zeros as well as infinitely many poles.
\end{proof}

This example has a pole at $s=1$ and does not satisfy the central-line-divisor hypothesis used above. It does show why nonnegative coefficients in \eqref{eq:Dpositive} cannot simply be called positivity of the difference datum. The needed extra information is the integer positive-definite \emph{coefficient function of the divisor} in Lemma~\ref{lem:PD}, followed by Cohen's theorem. The two uses of positivity have different mathematical content.

The coefficient sequences $f$ and $a_{\mathcal F}$ in \eqref{eq:bookerL} must also remain distinct. In particular, the third criterion in \eqref{eq:bookercriteria} concerns the ordinary Dirichlet coefficients, not $f(n)$. Exchanging them would make an invalid shortcut at precisely the critical exponent relevant to this problem.

\section{Four certified fields}\label{sec:certificate}

We now verify the real-zero condition for a member of the stipulated family. The explicit-formula approach to real zeros of five-dimensional $S_5$ Artin functions is established in \citet{dwyer2014}; the general analytic formula used below is in \citet[Section 3, equation (3--1)]{booker2006}. We give the test function, local data, and finite interval bounds explicitly. No claim of priority for the defining polynomial is made.

\begin{theorem}[A certified unconditional example]\label{thm:8249}
Let $L$ be the splitting field of
\[
f(x)=x^5-2x^4-x^3+4x+1.
\]
Then $\Gal(L/\Q)=S_5$, complex conjugation and $\operatorname{Frob}_5$ have cycle type $(12)(34)$, and $5$ is unramified. For $E$, $F$, and $H$ defined as above,
\[
\frac{\zeta_H(\sigma)}{\zeta_F(\sigma)}\ne0
\qquad(0<\sigma<1,\quad\sigma\in\R),
\]
where the quotient is understood as its entire Hecke continuation. Consequently all the conclusions of Theorem~\ref{thm:main} hold unconditionally for $L$. Both $L(s,\eta)$ and $L(s,\rh)$ have order zero at $s=1/2$.
\end{theorem}

The proof occupies this section. It has two independent mathematical parts: exact field and local-factor calculations, followed by a finite explicit-formula inequality. The latter certifies an interval of nonvanishing through a sum over all zeros; it does not infer interval nonvanishing from sampled values.

\subsection{Exact field data and ramified factors}

The discriminant is
\[
\operatorname{disc}(f)=8249=73\cdot113.
\]
Since it is squarefree, it is also the discriminant of the quintic field. An exact Sturm calculation gives one real root. For completeness, the leading coefficients of the Sturm chain, in descending degrees $5,4,3,2,1,0$, are
\[
1,\quad5,\quad\frac{26}{25},\quad-\frac{2450}{169},\quad
-\frac{30589}{480200},\quad\frac{1980584900}{5536609}.
\]
The variations at $-\infty$ and $+\infty$ are three and two, respectively. Thus the signature is $(1,2)$, giving the asserted complex conjugation.

The polynomial is irreducible modulo $2$ and has factorization degrees $(2,3)$ modulo $13$. Its group is therefore transitive of prime degree five and contains a transposition. To recall why this implies $S_5$, the graph whose edges are the conjugates of that transposition has components forming a block system. Primitivity makes the graph connected, and its edge transpositions generate $S_5$. At five the exact factorization is
\[
\overline f=(x+1)(x^2+3x+3)(x^2+4x+2).
\]
The two quadratics are irreducible and the factors are distinct. This proves the required unramified Frobenius type.

Put $\mathcal R=\eta\oplus\rh$. For a permutation $u\in S_5$, let $a(u)$ be its number of fixed points and $b(u)$ its number of two-cycles. The character needed for the logarithmic derivative is
\begin{equation}\label{eq:traceformula}
\chi_{\mathcal R}(u)=\operatorname{sgn}(u)
+\frac{a(u)(a(u)-3)}2+b(u).
\end{equation}
Indeed, the permutation representation on two-element subsets is
$\one\oplus\rho_4\oplus\rh$, whereas the permutation representation on points is $\one\oplus\rho_4$. Their difference has character
$\binom{a(u)}2+b(u)-a(u)$, as required. Thus ordinary finite-field factorization computes every unramified prime-power trace exactly.

The two ramified factorizations are
\begin{align*}
 f\bmod73&=(x+22)^2(x^3+27x^2+6x+27),\\
 f\bmod113&=(x+36)^2(x+90)(x^2+62x+33).
\end{align*}
The cubic and quadratic displayed here are irreducible over their respective finite fields. At either prime the inertia group is generated by a transposition: ramification is tame, and the discriminant exponent is one. The character $\eta$ has no inertia invariants. On $\rh^I$, the permutations of the three remaining letters have character values $(3,1,0)$ on the classes $1,2,3$, so this is the three-point permutation representation of $S_3$. This assertion can also be checked directly from
\[
\operatorname{tr}(u\mid\rh^I)
=\tfrac12\bigl(\chi_{\rh}(u)+\chi_{\rh}(u\tau)\bigr),
\qquad I=\langle\tau\rangle.
\]
The ramified Euler factors are consequently
\begin{equation}\label{eq:bad8249}
 B_{73}(s)=\frac1{1-73^{-3s}},\qquad
 B_{113}(s)=\frac1{(1-113^{-s})^2(1+113^{-s})}.
\end{equation}
In particular, the ramified traces are $t(73,k)=3\,\one_{3\mid k}$ and $t(113,k)=2+(-1)^k$.

The dimension of $\mathcal R^I$ is three, so its conductor exponent is three at both ramified primes. At infinity the trace of a double transposition on $\mathcal R$ is two; hence the $+1$ and $-1$ multiplicities are four and two. The conductor and completed function are therefore
\begin{equation}\label{eq:completion8249}
 Q_B=8249^3=561311462249,\qquad
 \Lambda_B(s)=Q_B^{s/2}\Gr(s)^4\Gr(s+1)^2B(s).
\end{equation}
The completion is entire because $B$ is the nontrivial quadratic Hecke function in \eqref{eq:B}. Its functional equation has sign $+1$:
$\Lambda_B(s)=\Lambda_B(1-s)$, as follows by taking the quotient of the completed Dedekind zeta functions. Its zeros lie in $0\le\Rea s\le1$ and satisfy the usual bound $N(T)=O(T\log(T+2))$. The completed endpoint values are nonzero: at $1$ the quotient is the ratio of the nonzero Dedekind residues, and at $0$ the functional equation applies.

\subsection{A test with nonnegative real part throughout the strip}

For $A>0$, define
\begin{equation}\label{eq:test8249}
 f_A(x)=\left(1-\frac{|x|}{A}\right)_+,\qquad
 g_A(x)=\frac{f_A(x)}{\cosh(x/2)},\qquad
 H_A(z)=\int_{\R}g_A(x)e^{izx}\dd x.
\end{equation}
The division by $\cosh(x/2)$ is the strip-positivity construction of \citet[Lemma 3.1]{miller}. We include the direct argument for this test.

\begin{lemma}[Strip positivity and the real-zero threshold]\label{lem:stripcertificate}
One has $\Rea H_A(z)\ge0$ for $|\Ima z|\le1/2$ and
$H_A(iy)\ge H_A(0)>0$ for real $y$. Suppose an entire completion of the form \eqref{eq:completion8249} is real on the real axis, has functional equation of sign $+1$, and has its zeros in the critical strip with counting function $O(T\log(T+2))$. If
\begin{equation}\label{eq:threshold}
\mathcal S_A:=\sum_{\Lambda_B(\rho)=0}
\Rea H_A\left(\frac{\rho-1/2}{i}\right)<2H_A(0),
\end{equation}
with zeros counted with multiplicity, it has no real zero in $(0,1)$, including at the center.
\end{lemma}

\begin{proof}
Evenness gives
\[
\Rea H_A(t\pm i/2)=\int_\R f_A(x)\cos(tx)\dd x
=A\left(\frac{\sin(At/2)}{At/2}\right)^2\ge0,
\]
with the continuous value at $t=0$. The function $g_A$ is continuous and its piecewise smooth derivative has bounded variation. Integration by parts twice, in the distributional sense at its three corners, gives
$H_A(t+iy)=O((1+|t|)^{-2})$ uniformly for $|y|\le1/2$. Apply the minimum principle to $\Rea H_A$ on rectangles bounded by $y=\pm1/2$ and $t=\pm T$. The lower bound on the vertical edges tends to zero as $T\to\infty$; hence the real part is nonnegative throughout the strip.

For real $y$,
\[
H_A(iy)=2\int_0^A g_A(x)\cosh(yx)\dd x\ge H_A(0)>0.
\]
The decay just proved and the zero-counting bound imply absolute convergence of the sum in \eqref{eq:threshold}. A real zero away from $1/2$ is accompanied by its distinct reflection and their total contribution is at least $2H_A(0)$. A central zero has even multiplicity because of the sign-$+1$ functional equation, so its contribution has the same lower bound. Every other zero has nonnegative contribution. Either kind of real zero contradicts \eqref{eq:threshold}.
\end{proof}

The assertion uses the full strip and therefore requires no Riemann-hypothesis assumption. In particular, unknown nonreal zeros cannot invalidate the lower bound contributed by a real zero.

\subsection{Explicit formula and a finite certificate}

Write
\[
-\frac{B'}B(s)=\sum_p\sum_{k\ge1}t(p,k)(\log p)p^{-ks}.
\]
For unramified $p$, the trace $t(p,k)$ is given by \eqref{eq:traceformula} at the $k$th power of its Frobenius permutation; at the ramified primes use \eqref{eq:bad8249}. In the normalization of \citet[Section 3, equation (3--1)]{booker2006} and \citet[Definition 1.3 and Examples 1.4]{booker2015}, the explicit formula is
\begin{equation}\label{eq:EF8249}
\mathcal S_A=C+2I_A-P_A,
\end{equation}
where
\begin{align}
 C&=\log Q_B-6\gamma-\pi-18\log2-6\log\pi,\label{eq:constant8249}\\
 K(x)&=\frac{4e^{-x/2}+2e^{-3x/2}}{1-e^{-2x}},\qquad
 I_A=\int_0^\infty K(x)(1-g_A(x))\dd x,\label{eq:gamma8249}\\
 P_A&=4\sum_{p^k\le e^A}\frac{t(p,k)\log p}{p^k+1}
 \left(1-\frac{k\log p}{A}\right).\label{eq:prime8249}
\end{align}
Here $\gamma$ is Euler's constant. The gamma shifts are $0,0,0,0,1,1$. Substitution of
\[
\psi(1/4)=-\gamma-\pi/2-3\log2,\qquad
\psi(3/4)=-\gamma+\pi/2-3\log2
\]
into the archimedean constant gives \eqref{eq:constant8249}. The prime coefficient is $4$, because the explicit formula contributes a factor $2$ and
$e^{-x/2}/\cosh(x/2)=2/(e^x+1)$. Since $g_A$ has compact support, \eqref{eq:prime8249} has no uncomputed prime tail.

The piecewise smooth test in \eqref{eq:test8249} is admissible by approximation. Convolve its triangular numerator with an even smooth approximate identity formed as a self-convolution, and then divide by $\cosh(x/2)$. The smooth explicit formula applies to these approximants. Their derivatives have uniformly bounded variation on a fixed compact interval, giving a uniform $O((1+|t|)^{-2})$ transform bound in the strip. Dominated convergence therefore applies to the absolutely convergent zero sum, the gamma integral, and the finite prime sums. This proves \eqref{eq:EF8249} with the stated test.

We take $A=8$. Directed interval evaluation proves $2980<e^8<2981$, so exactly $465$ prime powers occur. Every trace in the finite sum is an exact integer. An explicit algorithm for obtaining them starts with a prime sieve up to $2980$ and computes the degrees of the irreducible factors of $f\bmod p$ from
$\gcd(f,X^{p^d}-X)$ for $1\le d\le5$. If the factor degrees give cycle multiplicities $m_d$, then
\[
a(u^k)=\sum_{d\mid k}dm_d,\qquad
b(u^k)=\frac{a(u^{2k})-a(u^k)}2,\qquad
\operatorname{sgn}(u^k)=(-1)^{k(5-\sum_d m_d)}.
\]
These formulas and \eqref{eq:traceformula} determine the unramified terms; the two explicit ramified factors determine the rest.

For a direct enclosure of the gamma integral set
$t=e^{-x/2}$ and $R(x)=x/(e^x-1)$, with $R(0)=1$. On $[0,A]$, the integrand becomes the nonsingular expression
\begin{equation}\label{eq:regularJ}
\begin{split}
J_A(x):=K(x)(1-g_A(x))
={}&\frac{2t(2+t^2)(1-t)}{(1+t)(1+t^2)^2}\\
&+\frac4A\frac{2+t^2}{(1+t^2)^2}R(x).
\end{split}
\end{equation}
Both $t$ and $R$ are decreasing for $x\ge0$. The latter assertion follows by differentiating: its numerator is $e^x(1-x)-1<0$ for $x>0$. Divide $[0,8]$ into $1024$ intervals of length $1/128$, enclose $t$ and $R$ by endpoint bounds on each interval, and evaluate \eqref{eq:regularJ} by interval arithmetic. Summing the resulting rectangles encloses the finite gamma integral. The remaining tail satisfies
\begin{equation}\label{eq:gammatail}
T_0\le\int_A^\infty K(x)\dd x\le\frac{T_0}{1-e^{-2A}},
\qquad
T_0=8e^{-A/2}+\frac43e^{-3A/2}.
\end{equation}
These inequalities follow by replacing $(1-e^{-2x})^{-1}$ by its lower bound $1$ and upper bound $(1-e^{-2A})^{-1}$. Finally, $g_A$ is decreasing on $[0,A]$, so left and right rectangle sums enclose
$2H_A(0)=4\int_0^A g_A(x)\dd x$.

\begin{table}[ht]
\centering
\begin{tabular}{@{}lc@{}}
\toprule
Quantity&Certified rational enclosure\\
\midrule
$C$&$(1.10362,\ 1.10363)$\\
$I_8$&$(3.7324,\ 3.7646)$\\
$P_8$&$(0.52675,\ 0.52676)$\\
$\mathcal S_8=C+2I_8-P_8$&$(8.04,\ 8.11)$\\
$2H_8(0)$&$(8.96,\ 9)$\\
\bottomrule
\end{tabular}
\caption{Exact terminating rational endpoints enclosing the interval computation.}
\label{tab:certificate8249}
\end{table}

Table~\ref{tab:certificate8249} is obtained with $160$-bit MPFR arithmetic, using explicit downward and upward rounding for every endpoint operation; the correctness guarantees are specified in \citet[Sections 4.4 and 5.7]{mpfr}. Polynomial operations, Sturm chains, primes, and traces use exact integer or rational arithmetic. The displayed comparisons are asserted against exact rational numbers, not against rounded decimal output. A separately written check uses exact finite-field factorization and counts the permutation traces on points and two-element subsets directly, with $192$-bit Arb arithmetic and $2048$ cells; it also proves the sufficient inequalities $\mathcal S_8<17/2$ and $2H_8(0)>89/10$.

In particular, the first computation proves
\[
\mathcal S_8<\frac{811}{100}<\frac{224}{25}<2H_8(0).
\]
Lemma~\ref{lem:stripcertificate} proves nonvanishing for every real $0<\sigma<1$. The exact local checks permit application of Theorem~\ref{thm:main}. Its holomorphy conclusion and the factorization $B=L(\eta)L(\rh)$ then show that both factors have order zero at the center. This completes the proof of Theorem~\ref{thm:8249}.

\subsection{A second certified field}

\begin{corollary}\label{cor:60061}
Let $L$ be the splitting field of
\[
f(x)=x^5+2x^3+2x^2-2x+1.
\]
Then the conclusions of Theorem~\ref{thm:8249} hold for this field as well.
\end{corollary}

\begin{proof}
The discriminant is $60061=17\cdot3533$, again squarefree. An exact Sturm calculation gives one real root; irreducibility modulo $3$ and factorization degrees $(2,3)$ modulo $19$ prove that the group is $S_5$. At five,
\[
f\bmod5=(x+2)(x^2+x+1)(x^2+2x+3),
\]
with the required distinct irreducible factors. The ramified factorizations have double linear factors $x+10$ at $17$ and $x+1769$ at $3533$. After removing those squared factors, the remaining factor degrees are $(1,2)$ and $(1,1,1)$, respectively. The inertia calculation already proved above therefore gives
\[
B_{17}(s)=\frac1{(1-17^{-s})^2(1+17^{-s})},\qquad
B_{3533}(s)=\frac1{(1-3533^{-s})^3}.
\]
Thus $Q_B=60061^3=216659470006981$ and the gamma factors remain $\Gr(s)^4\Gr(s+1)^2$. The ramified prime $3533$ lies outside the finite prime-sum cutoff, but its conductor contribution is included.

Use the same test with $A=8$. Exact trace calculation, $160$-bit MPFR arithmetic, and $2048$ interval cells give enclosures
\[
8.8304<\mathcal S_8<8.8626,\qquad
8.9679<2H_8(0)<8.9836.
\]
The asserted rational comparisons are
\[
\mathcal S_8<\frac{89}{10}<\frac{224}{25}<2H_8(0).
\]
A separate exact-factorization and $192$-bit Arb calculation proves the same comparisons. Lemma~\ref{lem:stripcertificate} excludes every real zero in $(0,1)$, and Theorem~\ref{thm:main} gives the automorphy and central-order assertions.
\end{proof}

\subsection{An autocorrelation certificate for discriminant 30301}

\begin{corollary}\label{cor:30301}
Let $L$ be the splitting field of $f(x)=x^5-2x^3+2x^2+1$. Then the conclusions of Theorem~\ref{thm:8249} hold for this field as well.
\end{corollary}

\begin{proof}
The polynomial and field discriminants are $30301=157\cdot193$, and an exact Sturm calculation gives signature $(1,2)$. The polynomial is irreducible modulo $3$, has factorization degrees $(2,3)$ modulo $23$, and satisfies
\[
f\bmod5=(x+3)(x^2+3)(x^2+2x+4).
\]
Thus its group is $S_5$ and complex conjugation has the required double-transposition type. At each ramified prime there is one double linear factor. The residual factor degrees are $(3)$ at $157$ and $(1,2)$ at $193$. The same tame inertia calculation gives
\[
B_{157}(s)=\frac1{1-157^{-3s}},\qquad
B_{193}(s)=\frac1{(1-193^{-s})^2(1+193^{-s})}.
\]
Consequently $Q_B=30301^3=27820881360901$, with the same four even and two odd gamma shifts.

Let the first sixteen weights be
\[
\begin{split}
(w_0,\ldots,w_{15})={}&(305,396,476,559,619,667,717,775,\\
&825,870,904,936,961,979,992,1000),
\end{split}
\]
and set $w_{31-j}=w_j$ for $0\le j\le15$. Define $u(t)=w_j$ on $[j/2,(j+1)/2)$ for $0\le j\le31$, and zero elsewhere. Put
\begin{equation}\label{eq:autocorr30301}
C_m=\sum_{j=0}^{31-m}w_jw_{j+m},\qquad
h(x)=\frac{\int_\R u(t)u(t+x)\dd t}{\int_\R u(t)^2\dd t},
\qquad 0\le m\le32,
\end{equation}
where the empty sum is zero. Exact integer calculation gives
\[
C_0=19457050,\quad C_1=19322800,\quad C_{32}=0,
\qquad\min_{0\le m<32}(C_m-C_{m+1})=93025.
\]
For $x=(m+\theta)/2$, $0\le\theta\le1$, overlap lengths give
\begin{equation}\label{eq:linearh30301}
h(x)=\frac{(1-\theta)C_m+\theta C_{m+1}}{C_0}.
\end{equation}
Thus $h$ is even, nonnegative, continuous, decreasing on $[0,16]$, piecewise affine, and supported on $[-16,16]$. Its Fourier transform is
$\widehat h(t)=|\widehat u(t)|^2/\int_\R u(x)^2\dd x\ge0$.

Set $g(x)=h(x)/\cosh(x/2)$ and $H(z)=\int_\R g(x)e^{izx}\dd x$. Then
$\Rea H(t\pm i/2)=\widehat h(t)\ge0$. The derivative of $g$ has bounded variation, so $H$ has uniform quadratic decay in the closed strip. The harmonic minimum principle used in Lemma~\ref{lem:stripcertificate} gives $\Rea H\ge0$ throughout that strip. Since $g\ge0$, one also has $H(iy)\ge H(0)>0$. The real-zero threshold is therefore again $2H(0)$, including for a central zero.

The explicit formula is now
\begin{equation}\label{eq:EF30301}
\begin{split}
\mathcal S_h&=C+2I_h-P_h,\qquad I_h=\int_0^\infty K(x)(1-g(x))\dd x,\\
P_h&=4\sum_{p^k\le e^{16}}\frac{t(p,k)\log p}{p^k+1}h(k\log p),
\end{split}
\end{equation}
where $C$ is \eqref{eq:constant8249} with the present conductor and $K$ is \eqref{eq:gamma8249}. The same bounded-variation approximation proves admissibility. A directed enclosure gives $8886110<e^{16}<8886111$. There are $595877$ prime-power terms, with exact traces obtained from the fixed-point and two-cycle formula and the two ramified factors.

For the integral write $t=e^{-x/2}$, $R(x)=x/(e^x-1)$, and $L_h(x)=(1-h(x))/x$, extended by $L_h(0)=2(C_0-C_1)/C_0$. On $[0,16]$,
\begin{equation}\label{eq:Jh30301}
K(x)(1-g(x))=
\frac{2t(2+t^2)(1-t)}{(1+t)(1+t^2)^2}
+4\frac{2+t^2}{(1+t^2)^2}R(x)L_h(x).
\end{equation}
On each half-unit segment $h(x)=a+bx$, so $L_h(x)=(1-a)/x-b$ is monotone; on the first segment it is constant. A grid of $16384$ cells respects every breakpoint. Endpoint enclosures for $t$, $R$, and $L_h$, followed by interval evaluation of \eqref{eq:Jh30301}, bound the finite integral. Formula~\eqref{eq:gammatail}, with $A=16$, bounds the gamma tail. Monotonicity of $g$ gives rectangle bounds for $2H(0)=4\int_0^{16}g(x)\dd x$.

The $160$-bit MPFR computation proves the exact rational comparisons
\begin{equation}\label{eq:bounds30301}
\mathcal S_h<11<\frac{111}{10}<2H(0),\qquad
2H(0)-\mathcal S_h>\frac1{10}.
\end{equation}
The prime sum is evaluated directly from \eqref{eq:linearh30301}: on each segment its numerator is an integer linear combination of $\log p$ and $k(\log p)^2$, with endpoints rounded outwards. A separate $192$-bit Arb calculation with $32768$ cells proves the same comparisons. Exact finite-field factorizations in PARI and FLINT, with different trace evaluations, agree for all $595877$ terms; a fresh PARI stream was also regenerated for the MPFR check.

Full-strip positivity excludes every real zero in $(0,1)$. Theorem~\ref{thm:main} gives all-place standard-factor equality and entire automorphy, and nonvanishing of $B$ makes both central orders zero.
\end{proof}

\subsection{A fourth field with wild ramification at two}

\begin{corollary}\label{cor:48976}
Let $L$ be the splitting field of $f(x)=x^5-2x^3+x+2$. Then the conclusions of Theorem~\ref{thm:8249} hold for this field as well.
\end{corollary}

\begin{proof}
The polynomial discriminant is $48976=2^4\cdot3061$, where $3061$ is prime. An exact Sturm chain has successive leading coefficients
\[
1,\quad5,\quad4/5,\quad1,\quad-125,\quad3061/15625,
\]
so the signature is $(1,2)$. Irreducibility modulo three and factorization degrees $(2,3)$ modulo thirteen prove that the group is $S_5$. At five,
\[
f\bmod5=(x+3)(x^2+3x+4)(x^2+4x+1),
\]
with distinct irreducible quadratic factors. The arithmetic assumptions of Theorem~\ref{thm:main} therefore hold; this field also satisfies Calegari's original condition at five.

We first determine the wild local factor and conductor. Translation gives
\[
f(x+1)=x^5+5x^4+8x^3+4x^2+2\equiv x^4(x+1)\pmod2.
\]
Coprime Hensel factorization gives a linear factor with odd constant term and a monic quartic congruent to $x^4$. Their product has constant term two, so the quartic is Eisenstein. Its root generates the full ring of integers of the totally ramified quartic extension, by \citet[Chapter I, Section 6, Proposition 17 and its corollary]{serrelocal}. The two local factors have unit resultant. Hence the quintic power order is maximal at two; at $3061$, the discriminant valuation one already forces maximality. The field discriminant is therefore $48976$, and the sign character cuts out $\Q(\sqrt{3061})$. Its fundamental discriminant is $3061\equiv5\pmod8$, so $\eta$ is unramified at two with $\eta(\operatorname{Frob}_2)=-1$.

The decomposition group fixes the linear root and is transitive on the other four. Inertia is transitive on those four roots and, since $\eta$ is unramified, lies in $A_4$. Thus inertia is either $V_4$ or $A_4$. For lower ramification groups $I_i$, the Artin conductor formula and the permutation conductor--discriminant identity give
\begin{equation}\label{eq:conductor48976}
a_2(V)=\sum_{i\ge0}\frac{|I_i|}{|I_0|}
\operatorname{codim}V^{I_i},\qquad a_2(M_1)=4,
\end{equation}
where $M_1$ is the permutation representation on the five roots; see \citet[Chapter VI, Section 2, Proposition 2, Corollary $1'$, and Proposition 4]{serrelocal}. If $I_0=V_4$, then $I_1=V_4$ because the tame quotient has odd order. The first two terms in \eqref{eq:conductor48976} would be $3+3=6$, a contradiction. Therefore $I_0=A_4$ and $I_1=V_4$. Their contributions are already $3+(4/12)3=4$. Every later contribution must vanish; faithfulness of $M_1$ forces $I_2=1$.

Let $M_2$ be the permutation representation on unordered pairs. The character identity
\[
\mathcal R=\eta+M_2-M_1
\]
computes the required invariants. Under $A_4$, both $M_1$ and $M_2$ have two orbits. Under $V_4$, $M_1$ has two and $M_2$ has four: the six pairs within the moving four-point set form three orbits, and the four pairs containing the fixed root form one. Hence
\[
\dim\mathcal R^{A_4}=1,\qquad
\dim\mathcal R^{V_4}=3,\qquad
 a_2(\mathcal R)=5+\frac4{12}\,3=6.
\]
The decomposition group preserves each of the two $A_4$-orbits in both permutation representations, so their invariant representations are both trivial of dimension two. Their virtual difference vanishes, leaving $\eta$ on $\mathcal R^{A_4}$. Consequently
\begin{equation}\label{eq:wildfactor48976}
B_2(s)=\frac1{1+2^{-s}},\qquad t(2,k)=(-1)^k.
\end{equation}

At the other ramified prime, exact factorization gives
\[
f\bmod3061=(x+419)(x+1151)^2(x^2+340x+2630),
\]
with the quadratic irreducible. Tame transposition inertia gives the invariant degrees $(1,2)$, so
\[
B_{3061}(s)=\frac1{(1-3061^{-s})(1-3061^{-2s})},
\qquad Q_B=2^6\cdot3061^3=1835565822784.
\]
The gamma shifts remain $0,0,0,0,1,1$.

The same triangular test with $A=8$ now suffices. All $465$ prime-power traces are recalculated, using \eqref{eq:wildfactor48976} at two. The prime $3061$ lies beyond the sum cutoff but is included in $Q_B$. At $160$ bits with $2048$ interval cells, directed MPFR arithmetic proves
\begin{equation}\label{eq:bounds48976}
\mathcal S_8<\frac{61}{10}<\frac{224}{25}<2H_8(0),
\qquad 2H_8(0)-\mathcal S_8>\frac{14}{5}.
\end{equation}
A separate exact maximal-order computation, enumeration of the transitive subgroups of $A_4$, and permutation-trace calculation, followed by $192$-bit Arb arithmetic with $4096$ cells, proves the same inequalities. All traces also agree with a fresh PARI factorization calculation. Lemma~\ref{lem:stripcertificate} excludes every real zero, and Theorem~\ref{thm:main} gives the claimed automorphy and local equality.
\end{proof}

An additional exact check expresses the wild conductor through a degree-ten field. The pair sums have polynomial
\[
q(y)=y^{10}-6y^8+9y^6-22y^5-4y^4+16y^3-4,
\]
as verified by
$\operatorname{Res}_x(f(x),f(y-x))=2^5f(y/2)q(y)^2$.
Put $K_2=\Q[y]/(q)$. Its power-order discriminant is $2^{18}2963^2 3061^3$, its index in the maximal order is $16\cdot2963$, and its field discriminant is $2^{10}3061^3$. The identity $B=L(s,\eta)\zeta_{K_2}/\zeta_K$ therefore recovers $Q_B=3061\,\Delta_{K_2}/\Delta_K=2^6 3061^3$. The direct inertia proof above supplies the wild local input independently of this additional maximal-order calculation.

\subsection{Comparison with a published example}

For comparison, \citet[Theorems 5.1 and 5.3]{calegari} proves automorphy and records Booker's certified calculation for
\[
f_0(x)=x^5-x^3-x^2+x+1,\qquad\Delta_K=1609.
\]
The published certificate gives the lowest positive zero height as $1.624\ldots$, states that all zeros with imaginary part of magnitude less than $100$ lie on the critical line, and explicitly gives nonvanishing throughout the real critical interval. This polynomial is irreducible modulo five, so its Frobenius there is a $5$-cycle. Calegari handled its modularity input by a separate two-adic construction in his Section 5. Pilloni--Stroh's theorem now supplies the input without that auxiliary construction, and the published nonvanishing of the five-dimensional factor suffices by Remark~\ref{rem:weaker}. The four certificates above establish real nonvanishing of the entire Hecke quotient $B$ directly.

\section{Scope, dependencies, and remaining questions}

The theorem establishes a conditional automorphy result with a strictly weaker real-zero hypothesis than Calegari's stated Theorem 1.2. It does not prove \eqref{eq:RBC} for every field in the family. The second Hecke function restricts a possible comparison pole to a common-zero locus; it does not prove that locus empty. Nor does an absence of off-central \emph{real} zeros locate nonreal zeros: GRH for the individual Artin factors is an additional question. The central order identity is exact; in the four certified examples of discriminants $8249$, $60061$, $30301$, and $48976$, both orders on its right are zero. For other fields the identity alone does not determine either order.

The automorphic construction, cuspidality, temperedness, and finite-exception comparison are imported from Calegari. The proof presented here supplies the analytic passage that excludes an infinite exceptional set while allowing a central zero. The essential non-elementary harmonic-analysis input is Cohen's idempotent theorem; Bochner's theorem and the discrete-coset consequence are used with their correct topologies. Hadamard factorization uses the order-one analytic structure of completed $L$-functions. None of these inputs is replaced by a conjectural positivity principle.

The argument also explains why a general statement about positive Euler products would be insufficient. The proof simultaneously uses integral meromorphic orders, a reflection functional equation, a divisor confined to the central line, finite-order growth, a finite coset decomposition, and bounded local prime-power coefficients. Removing one of these requirements requires a separate theorem; the example in Proposition~\ref{prop:counterexample} already prevents one common generalization.

\subsection{The unresolved off-central real pole}

Without \eqref{eq:RBC}, Lemma~\ref{lem:abscissa} still supplies a real pole at an abscissa $\alpha\in(0,1)$ if $\Omega$ is infinite. Reflection of that pole rules out $\alpha<1/2$: the reflected pole would lie in the zero-free Euler half-plane. If $\alpha=1/2$, Sections 5--7 apply without any further use of \eqref{eq:RBC} and give the same contradiction. Thus the only residual possibility is
\[
\tfrac12<\alpha<1,\qquad B(\alpha)=C(\alpha)=0.
\]
The quotient then has a pole at both $\alpha$ and $1-\alpha$, and its divisor lies in the closed strip between these two lines.

Using \citet[Theorems 1.2 and 8.29]{liu2026} as an additional external input narrows this interval to
\[
\frac12<\alpha\le B_*,
\qquad
B_*=0.873345856106718576801184796065\ldots .
\]
Here \(B_*\) is the degree-nine algebraic constant defined by equations (42) and (44) of the cited paper. It is uniquely isolated, for the minimal polynomial printed there, by
\[
\frac{8733458561067185768}{10^{19}}
<B_*<
\frac{8733458561067185769}{10^{19}}.
\]
Liu's Theorem 1.2 gives a zero-free half-plane for every continuous unitary Hecke character over every number field; Theorem 8.29 is its finite-order and norm-twist case. Thus it applies to the nontrivial quadratic character of \(H/F\). A comparison pole at \(\alpha\) forces a zero of \(L(s,\rh)\) there and hence a zero of the entire function \(B=L(s,\eta)L(s,\rh)\); the cited half-plane excludes \(\alpha>B_*\). The simpler one-stage Theorem 1.1 gives the weaker rational endpoint \(3493459/4000000\). This refinement uses the named external zero-free theorem, whose full analytic proof is not reproduced or independently established here. Neither Theorem~\ref{thm:main} nor any of the four finite certificates depends on it.

The Abel argument still proves that
\[
r_\alpha(t)=-\ord_{s=\alpha+it}M(s)
=\lim_{\delta\downarrow0}\delta D_\Omega(\alpha+\delta+it)
\]
is integer-valued, positive definite on $\R_d$, bounded, and locally finitely supported. Cohen's theorem still gives a finite progression description of this \emph{boundary} divisor. It gives no description of the divisor strictly inside the strip. Accordingly a product constructed from $r_\alpha$ need not have the full divisor of $M$, and the zero-free Hadamard quotient in \eqref{eq:Hadamard} is unavailable.

There is also a precise obstruction to transferring the original prime support through an ordinary Fourier limit. Define positive finite measures
\[
\nu_\delta=4\delta\sum_{p\in\Omega}\sum_{\substack{k\ge1\\k\ \mathrm{odd}}}
(\log p)p^{-k(\alpha+\delta)}\delta_{k\log p}.
\]
Their total masses tend to $r_\alpha(0)>0$. On every fixed compact interval only finitely many prime powers occur, and each weight tends to zero. Thus these measures converge vaguely to zero but are not tight; their Fourier transforms have a pointwise limit on $\R_d$, not the Fourier transform of a narrow limit of finite measures on the ordinary real line.

The central prime measure
\[
\mu=4\sum_{p\in\Omega}\sum_{\substack{k\ge1\\k\ \mathrm{odd}}}
(\log p)p^{-k/2}\bigl(\delta_{k\log p}+\delta_{-k\log p}\bigr)
\]
cannot be tempered in this residual case. A positive tempered measure has polynomially bounded mass on intervals $[-T,T]$. Integrating $e^{-(\sigma-1/2)x}$ on its positive half would then converge for every $\sigma>1/2$, giving absolute convergence of $D_\Omega(\sigma)$ there and contradicting $\alpha>1/2$.

For a simple test of what the boundary information permits, $r(t)=m\one_{\{0\}}(t)$ is integer-valued and positive definite on $\R_d$ for every positive integer $m$. It has locally finite support, and its ordinary atomic measure has a constant continuous Fourier density. This is a countermodel to an inference from boundary positivity to pure atomic prime support; it is not an arithmetic example or a constructed comparison $L$-function. The off-central real-pole alternative remains unexcluded by the present argument. Appendix~\ref{sec:offcentral-obstructions} strengthens this observation: it proves the canonical removal obstruction for every multiplicity, gives a sufficient rigidity condition on the entire divisor, and constructs explicit examples showing why neither boundary information nor a presumed uniform interior gap supplies that condition.

\appendix
\section{An archimedean consistency check}

The gamma elimination in Lemma~\ref{lem:gamma} is robust enough to start with an unknown integer $d$. In this particular family one can also verify $d=0$ directly from the construction. The two-dimensional finite-image representation over $E$ used by Calegari is odd at both real places, with parameter $\one\oplus\operatorname{sgn}$. Its normalized fourth symmetric power has parameter
\[
\operatorname{Sym}^4(\one\oplus\operatorname{sgn})
\otimes\det(\one\oplus\operatorname{sgn})^{-2}
=\one^{\oplus3}\oplus\operatorname{sgn}^{\oplus2}.
\]
This is the archimedean parameter of the base change of $\pi$. Because $E$ is real quadratic, local base change at a real place is the identity. Hence $\pi_\infty$ has exactly the factor \eqref{eq:gamma}. This check uses local compatibility in Calegari's construction over $E$, rather than assuming the desired global comparison over $\Q$.

\section{A Fourier interpretation of the arithmetic obstruction}

The progression conclusion also has a distributional interpretation. If an integer-valued positive-definite function $r$ has locally finite support, Lemma~\ref{lem:cosets} gives a finite sum of lattice combs and point masses for
$\nu=\sum_t r(t)\delta_t$. With Fourier convention
$\widehat\nu(x)=\int e^{itx}\dd\nu(t)$, Poisson summation gives
\[
\widehat{\sum_{n\in\Z}\delta_{\beta+\alpha n}}
=\frac{2\pi}{\alpha}\sum_{m\in\Z}
e^{2\pi im\beta/\alpha}\delta_{2\pi m/\alpha}.
\]
An isolated atom transforms into a continuous exponential density. If an independently justified explicit formula makes $\widehat\nu$ purely atomic, the continuous part must vanish, and its support lies in finitely many frequency lattices. Each lattice can contain $\log p$ for at most one rational prime, by unique factorization. Nonzero prime atoms for infinitely many different primes are therefore impossible.

This explains the Cohen mechanism in ordinary Fourier language. A complete proof through this route must establish temperedness, the precise Fourier normalization, equality of the gamma kernels, and the treatment of finite bad-prime corrections. The main proof avoids requiring those distributional steps by reconstructing the meromorphic quotient directly. It nevertheless uses the same integral positive-definite divisor structure.

\section{Boundary divisors and canonical factor removal}\label{sec:offcentral-obstructions}

This appendix studies the boundary function \(r_\alpha(t)=-\ord_{\alpha+it}M\) when the comparison abscissa lies strictly to the right of \(1/2\). It proves the two assertions in Theorem~\ref{thm:boundary-method} and gives explicit examples delimiting their scope. Throughout, positive definiteness is on \(\R_d\), while local finiteness refers to the ordinary topology on \(\R\). A generalized Dirichlet series has distinct positive frequencies that are locally finite in \([0,\infty)\), and converges absolutely in some right half-plane.

The canonical factor of a boundary divisor means the finite product of exponential and linear factors in \eqref{eq:ob3-9}, normalized so that its logarithmic derivative has no constant term at \(+\infty\). The results concern this precise operation. They do not replace the nonvanishing condition \eqref{eq:RB}, and do not narrow the remaining interval for \(\alpha\).

\subsection{Simple boundary profiles}

\begin{proposition}[Simple boundary profiles]\label{prop:ob-simple}

Let \(r:\mathbb R_d\to\mathbb Z\) be positive definite, with \(r(0)=1\), and suppose its support is locally finite in the ordinary real topology. Exactly one of the following forms occurs:

\begin{equation}\label{eq:ob2-1}
r=\mathbf1_{\{0\}};
\end{equation}

\begin{equation}\label{eq:ob2-2}
r=\mathbf1_{h\mathbb Z}\quad\text{for some }h>0;
\end{equation}

\begin{equation}\label{eq:ob2-3}
r(nh)=(-1)^n\quad(n\in\mathbb Z),\qquad
r(t)=0\quad(t\notin h\mathbb Z),
\end{equation}

again for some \(h>0\). No idempotent theorem is needed for this special case.

\end{proposition}

\begin{proof} Positivity gives \(|r(t)|\le1\), so the values are \(-1,0,1\). Real-valuedness also gives \(r(-t)=r(t)\). Realize the kernel \(r(t-u)\) as a Gram kernel of unit vectors \(v_t\). This realization can be obtained directly by completing the space of finite formal linear combinations modulo its null space; it invokes no continuity of \(r\).

If \(|r(t)|=1\), equality in Cauchy--Schwarz gives \(v_t=r(t)v_0\). Translation is an isometry of the Gram construction, and therefore

\begin{equation}\label{eq:ob2-4}
r(t+u)=r(t)r(u)\qquad(u\in\mathbb R).
\end{equation}

Consequently \(H=\{t:|r(t)|=1\}\) is a subgroup and \(r|_H\) is a character with values in \(\{1,-1\}\). Since the other values of \(r\) are zero, \(H\) is its support. A locally finite subgroup of \(\mathbb R\) is either \(\{0\}\) or \(h\mathbb Z\): in the nonzero case its positive elements have a least member \(h\), and Euclidean division reduces every member modulo \(h\). A character of \(h\mathbb Z\) into \(\{1,-1\}\) is determined by its value at \(h\). These are precisely \eqref{eq:ob2-1}--\eqref{eq:ob2-3}. Conversely all three are positive definite, by grouping finite Gram sums according to cosets of \(H\). \end{proof}

\paragraph{Arithmetic interpretation.}

If the entire quadratic Hecke function \(B=L(s,\eta_E)L(s,\rho_5)\) has a \emph{simple} zero at a real comparison pole \(\alpha\), then \(m=1\). Indeed, \(\Lambda(s,\pi)\) is entire at \(\alpha\), so a pole of \(M=\Lambda(\pi)/\Lambda(\rho_5)\) forces a positive zero order of \(L(s,\rho_5)\). The entire factor \(L(s,\eta_E)\) has nonnegative zero order there. Their sum is one, so \(L(s,\rho_5)\) has order one, \(L(s,\eta_E)\) has order zero, and the numerator has order zero. This implication is conditional on simplicity; simplicity of an arbitrary off-central zero of \(B\) is not asserted.

\subsection{Canonical boundary removal}

Here ``canonical removal'' means dividing by the finite-order elementary factor having exactly the boundary divisor in Proposition~\ref{prop:ob-simple}, with no extra divisors. The harmless freedom to multiply that factor by \(e^{as+b}\) only changes a conductor constant in its logarithmic derivative. We use the Euler normalization in which that constant is zero. Arbitrary zero-free analytic factors are not permitted: allowing them would simply move the unknown part of \(Q_\Omega\) into the factor being removed.

\begin{proposition}[Removal of a simple boundary divisor]\label{prop:ob-simple-peeling}

Suppose \(\alpha>0\), and \(D_\Omega\) is as in \eqref{eq:Dpositive}, absolutely convergent on a right half-plane. For each of the three simple-pole profiles in Proposition~\ref{prop:ob-simple}, let \(P_\alpha\) be the canonical factor below. Then

\begin{equation}\label{eq:ob3-1}
-\frac{(Q_\Omega/P_\alpha)'}{Q_\Omega/P_\alpha}
=D_\Omega-D_{P_\alpha}
\end{equation}

cannot be an absolutely convergent generalized Dirichlet series with nonnegative coefficients and positive frequencies. The same conclusion holds if the bounded finite bad-prime contributions in \eqref{eq:Dar} are included.

\end{proposition}

\begin{proof}\textit{Singleton profile.} Take

\begin{equation}\label{eq:ob3-2}
P_\alpha(s)=(s-\alpha)^{-1},\qquad
D_{P_\alpha}(s)=\frac1{s-\alpha}.
\end{equation}

Any absolutely convergent prime-power series in \eqref{eq:Dpositive} or \eqref{eq:Dar} is \(O(e^{-c\sigma})\) as real \(\sigma\to+\infty\), for some \(c>0\). Therefore

\begin{equation}\label{eq:ob3-3}
D_\Omega(\sigma)-\frac1{\sigma-\alpha}<0
\end{equation}

for all sufficiently large \(\sigma\). This is incompatible with a nonnegative Dirichlet series. In Laplace-transform language, \(1/(s-\alpha)\) corresponds to the continuous positive density \(e^{\alpha x}\,dx\) on \(x>0\), whereas \(D_\Omega\) comes from a positive atomic measure. Its removal subtracts that continuous density. A conductor constant at frequency zero does not remove the negative continuous part.

\textit{Lattice profile.} If \(r=\mathbf1_{h\mathbb Z}\), put \(\lambda=2\pi/h\) and

\begin{equation}\label{eq:ob3-4}
P_\alpha(s)=(1-e^{-\lambda(s-\alpha)})^{-1}.
\end{equation}

Then, for \(\Re s>\alpha\),

\begin{equation}\label{eq:ob3-5}
D_{P_\alpha}(s)=\lambda\sum_{n\ge1}
e^{\alpha n\lambda}e^{-n\lambda s}.
\end{equation}

At frequencies \(n\lambda\), the prime-power series can involve at most one underlying prime. In fact, \(n\lambda=k\log p\) and \(n'\lambda=k'\log q\) imply \(p^{kn'}=q^{k'n}\), hence \(p=q\). The original coefficient on this entire frequency progression is consequently bounded by \(4\log p\), or by \(10\log p\) when finite bad-prime terms are included; it is zero when no prime occurs. The subtracted coefficient \(\lambda e^{\alpha n\lambda}\) tends to infinity. Thus the remainder has negative coefficients at arbitrarily large \(n\lambda\).

\textit{Alternating lattice profile.} If \(r(nh)=(-1)^n\), put \(\lambda=\pi/h\) and

\begin{equation}\label{eq:ob3-6}
P_\alpha(s)=
\frac{1+e^{-\lambda(s-\alpha)}}{1-e^{-\lambda(s-\alpha)}}.
\end{equation}

Its poles are at the even multiples of \(ih\), and its zeros at the odd multiples, giving exactly the required signed divisor. Moreover

\begin{equation}\label{eq:ob3-7}
D_{P_\alpha}(s)=2\lambda
\sum_{\substack{n\ge1\\n\text{ odd}}}
e^{\alpha n\lambda}e^{-n\lambda s}.
\end{equation}

The same bounded-coefficient argument gives negative remainder coefficients at all sufficiently large odd model frequencies that occur in the union of frequency sets.

In the last two cases the union of the lattice frequencies and prime-power logarithms is locally finite and bounded away from zero. Both series converge absolutely sufficiently far to the right. Uniqueness of generalized Dirichlet series applies: if a nonzero coefficient at the least remaining frequency existed in a series representing zero, multiplying by its leading exponential and letting \(\sigma\to+\infty\) would contradict that coefficient. Iterating proves uniqueness. Thus the negative coefficients cannot be replaced by another nonnegative expansion. \end{proof}

The argument also survives canonical removal of the reflected boundary at \(1-\alpha>0\). The reflected divisor is the same in the three cases because the profiles are even. The corresponding positive-frequency logarithmic derivatives are again nonnegative, with \(\alpha\) replaced by \(1-\alpha\), so their subtraction cannot correct the negative terms just exhibited.

\paragraph{Scope of the obstruction.} Proposition~\ref{prop:ob-simple-peeling} excludes an iterative Landau/Cohen argument that requires the canonically peeled remainder to retain a nonnegative Euler logarithmic derivative. It does not exclude a different invariant, a signed explicit formula, or an argument that uses the complete automorphic and Artin structure simultaneously.

\begin{proposition}[Canonical removal at arbitrary multiplicity]\label{prop:ob-peeling}

The conclusion of Proposition~\ref{prop:ob-simple-peeling} holds for every nonzero integer-valued positive-definite boundary profile with locally finite ordinary support, without assuming \(r(0)=1\). More precisely, start with any Cohen decomposition

\begin{equation}\label{eq:ob3-8}
r(t)=\sum_{j=1}^J c_j\mathbf1_{\beta_j+h_j\mathbb Z}(t)
+\sum_{u=1}^U b_u\mathbf1_{\{t_u\}}(t),
\qquad c_j,b_u\in\mathbb Z,\quad h_j>0.
\end{equation}

At a line of positive real part \(\alpha\), its canonical factor is

\begin{equation}\label{eq:ob3-9}
P_\alpha(s)=
\prod_j(1-e^{-\lambda_j(s-\alpha-i\beta_j)})^{-c_j}
\prod_u(s-\alpha-it_u)^{-b_u},
\qquad \lambda_j=2\pi/h_j.
\end{equation}

Then \(D_{P_\alpha}=-P_\alpha'/P_\alpha\) is the Laplace transform of a nonzero nonnegative measure on positive frequencies. Nevertheless \(D_{\rm ar}-D_{P_\alpha}\) cannot have a nonnegative generalized Dirichlet expansion when \(D_{\rm ar}\) satisfies \eqref{eq:Dar} and \eqref{eq:bound}. This is a statement about positivity of the \emph{logarithmic derivative}, not a claim that every possible ordinary coefficient of the quotient is negative.

\end{proposition}

\begin{proof}\textit{Positivity of the canonical factor.} Let \(K=\widehat{\mathbb R_d}\) be the compact group of all characters of the additive discrete real group. Adopt the convention

\begin{equation}\label{eq:ob3-10}
r(t)=\int_K\chi(t)\,d\nu(\chi).
\end{equation}

Bochner's theorem for locally compact abelian groups \citep[Theorem~12.3]{korner2008} supplies a positive finite measure \(\nu\). Let \(m_K\) be Haar probability measure on \(K\). For \(H_j=h_j\mathbb Z\) put

\begin{equation}\label{eq:ob3-11}
K_j=H_j^\perp=\{\chi\in K:\chi(h_j)=1\},
\end{equation}

and let \(m_j\) be Haar probability measure on \(K_j\). Orthogonality gives

\begin{equation}\label{eq:ob3-12}
\widehat{\overline{\chi(\beta_j)}\,m_j}(t)
=\mathbf1_{\beta_j+H_j}(t),\qquad
\widehat{\overline{\chi(t_u)}\,m_K}(t)
=\mathbf1_{\{t_u\}}(t).
\end{equation}

Fourier-transform uniqueness therefore identifies the positive measure in \eqref{eq:ob3-10} as

\begin{equation}\label{eq:ob3-13}
\nu=F_0m_K+\sum_{j=1}^J
c_j\overline{\chi(\beta_j)}\,m_j,
\qquad
F_0(\chi)=\sum_u b_u\overline{\chi(t_u)}.
\end{equation}

Partition the \(h_j\) by rational commensurability. For a class \(I\), choose \(h_I>0\) such that

\begin{equation}\label{eq:ob3-14}
d_j:=h_I/h_j\in\mathbb Z_{>0}\quad(j\in I).
\end{equation}

Then \(H_I=h_I\mathbb Z\) is contained in every \(H_j\) in the class with finite index \([H_j:H_I]=d_j\). Its annihilator \(K_I\) contains each \(K_j\) as a clopen subgroup of index \(d_j\), and

\begin{equation}\label{eq:ob3-15}
m_j=d_j\mathbf1_{K_j}\,m_I.
\end{equation}

The part of \(\nu\) belonging to this class is thus

\begin{equation}\label{eq:ob3-16}
\nu_I=F_I m_I,\qquad
F_I(\chi)=\sum_{j\in I}
c_jd_j\,\overline{\chi(\beta_j)}\,\mathbf1_{K_j}(\chi),
\end{equation}

where \(F_I\) is continuous on \(K_I\).

The measures \(m_K\) and the finitely many \(m_I\) are mutually singular. Indeed, \(m_K(K_I)=0\): the evaluation map \(\chi\mapsto\chi(h_I)\) maps Haar measure to Haar measure on the circle, since \(h_I\) has infinite order in \(\mathbb R_d\). For distinct classes \(I,I'\), the class of \(h_{I'}\) has infinite order in \(\mathbb R_d/H_I\). The same orthogonality argument gives

\begin{equation}\label{eq:ob3-17}
m_I(K_I\cap K_{I'})=0
\quad\text{and}\quad
m_{I'}(K_I\cap K_{I'})=0.
\end{equation}

For example, the pushforward under \(\chi\mapsto\chi(h_{I'})\) has all nonzero Fourier coefficients zero because \(nh_{I'}\notin H_I\) for \(n\ne0\), and hence is circle Haar measure. This proves \eqref{eq:ob3-17} without assuming any metrizability of \(K\).

Since \(\nu\) is positive and these components are mutually singular, \(F_0m_K\) and every \(F_Im_I\) are positive. Haar measure has full support on each compact group, so continuity upgrades almost-everywhere nonnegativity to

\begin{equation}\label{eq:ob3-18}
F_0(\chi)\ge0\quad(\chi\in K),\qquad
F_I(\chi)\ge0\quad(\chi\in K_I).
\end{equation}

Evaluate first at the ordinary characters \(\chi_x(t)=e^{-ixt}\). This gives the nonnegative trigonometric polynomial

\begin{equation}\label{eq:ob3-19}
f_0(x):=F_0(\chi_x)=\sum_u b_u e^{it_ux}\ge0.
\end{equation}

For class \(I\), put \(\lambda_I=2\pi/h_I\). At \(x=n\lambda_I\), the character \(\chi_x\) lies in \(K_I\), and it belongs to \(K_j\) precisely when \(d_j\mid n\). The coefficient of \(e^{-n\lambda_I s}\) in the progression part of \(D_{P_\alpha}\) is therefore

\begin{equation}\label{eq:ob3-20}
e^{\alpha n\lambda_I}u_I(n),\qquad
u_I(n)=\lambda_I F_I(\chi_{n\lambda_I})\ge0.
\end{equation}

The function \(u_I(n)\) is a finite sum of unit-modulus exponentials, because each divisibility indicator is a finite root-of-unity average. Distinct classes share no positive frequency. The partial fractions from the individual points have the integral representation

\begin{equation}\label{eq:ob3-21}
\sum_u\frac{b_u}{s-\alpha-it_u}
=\int_0^\infty e^{-sx}e^{\alpha x}f_0(x)\,dx.
\end{equation}

Equations \eqref{eq:ob3-19}--\eqref{eq:ob3-21} show that

\begin{equation}\label{eq:ob3-22}
D_{P_\alpha}(s)=
\int_0^\infty e^{-sx}e^{\alpha x}f_0(x)\,dx
+\sum_I\sum_{n\ge1}
e^{\alpha n\lambda_I}u_I(n)e^{-n\lambda_I s}
\end{equation}

is a nonnegative Laplace transform for \(\Re s>\alpha\).

\textit{Nontriviality.} If \(f_0\) vanishes identically, linear independence of the distinct real exponentials gives \(F_0=0\). If \(u_I(n)=0\) for every positive integer \(n\), the finite-exponential-sum mean-square argument gives that its coefficients vanish, so it also vanishes for all integers \(n\). The ordinary characters \(\{\chi_{n\lambda_I}:n\in\mathbb Z\}\) are dense in \(K_I\): their annihilator is exactly \(h_I\mathbb Z\), and the double-annihilator theorem \citep[Lemma~12.5.2]{dikranjan2018} identifies their closure. Continuity then gives \(F_I=0\). Thus if every term on the right of \eqref{eq:ob3-22} vanished, \eqref{eq:ob3-13} would give \(\nu=0\), hence \(r=0\). The profile is assumed nonzero, so at least one part is nonzero.

\textit{Failure of positivity after removal.} If \(f_0\ne0\), its nonnegative real-analytic function on the real axis has a first nonzero Taylor term at zero, of even degree \(2q\), with a positive coefficient \(c\). Taylor expansion on a small interval, followed by an exponentially bounded tail, gives

\begin{equation}\label{eq:ob3-23}
\int_0^\infty e^{-(\sigma-\alpha)x}f_0(x)\,dx
\sim \frac{c(2q)!}{(\sigma-\alpha)^{2q+1}}
\quad(\sigma\to+\infty).
\end{equation}

The other terms in \eqref{eq:ob3-22} are nonnegative. In contrast, \(D_{\rm ar}(\sigma)=O(e^{-c_1\sigma})\) for some \(c_1>0\). Therefore the real part of \(D_{\rm ar}(\sigma)-D_{P_\alpha}(\sigma)\) is eventually negative. It cannot be a nonnegative Dirichlet series.

If \(f_0=0\), choose a nonzero class \(I\). A nonzero finite exponential sum \(u_I(n)\ge0\) is at least some \(\epsilon>0\) along an infinite sequence of positive integers: otherwise it would tend to zero, contradicting its positive Cesàro mean square. At these frequencies the subtracted coefficients

\begin{equation}\label{eq:ob3-24}
e^{\alpha n\lambda_I}u_I(n)
\end{equation}

tend to infinity, whereas the coefficient of \(D_{\rm ar}\) on this progression is uniformly bounded by the one-prime argument in Proposition~\ref{prop:ob-simple-peeling}. Hence infinitely many remainder coefficients have negative real part. Uniqueness of generalized Dirichlet series finishes the proof. \end{proof}

The argument also applies after removing both reflected boundaries, because each has a positive canonical logarithmic derivative. It supplies no replacement positivity for the interior divisor. It is a rigorous obstruction to this particular induction scheme, not a proof of the impossibility of other off-central reductions.

\subsection{Rigidity from the complete divisor}

The missing input is not the numerical value \(1/2\) in the coefficient comparison. It is control of the \emph{entire} divisor.

\begin{proposition}[Rigidity from finitely many complete progressions]\label{prop:ob-global}

Assume the comparison representation \eqref{eq:M}, the gamma and reflection equations \eqref{eq:G} and \eqref{eq:FE}, and the coefficient bound \eqref{eq:bound}. Suppose the entire signed divisor \(-\operatorname{ord}M\) is a finite integral combination of vertical arithmetic progressions and finitely many individual points, all lying in \(0<\Re s<1\). Then \(M\) is constant and \(D_{\rm ar}=0\). In particular \(\Omega\) is empty under this exact Euler representation.

\end{proposition}

\begin{proof} Write the progression terms as

\begin{equation}\label{eq:ob4-1}
\sum_{j=1}^J c_j\mathbf1_{
\{\sigma_j+i(\beta_j+h_jn):n\in\mathbb Z\}},
\qquad 0<\sigma_j<1,
\end{equation}

and the point terms as \(\sum_u b_u\mathbf1_{\{\rho_u\}}\), where \(0<\Re\rho_u<1\). Set \(\lambda_j=2\pi/h_j\) and

\begin{equation}\label{eq:ob4-2}
P(s)=\prod_j
(1-e^{-\lambda_j(s-\sigma_j-i\beta_j)})^{-c_j}
\prod_u(s-\rho_u)^{-b_u}.
\end{equation}

This is a meromorphic function of order at most one with the entire divisor of \(M\). Hadamard factorization therefore gives \(M=e^{as+b}P\). In particular

\begin{equation}\label{eq:ob4-3}
-M'/M=-a+R+C,
\quad R(s)=\sum_u\frac{b_u}{s-\rho_u},
\end{equation}

\begin{equation}\label{eq:ob4-4}
C(s)=\sum_j c_j\lambda_j\sum_{n\ge1}
e^{\sigma_j n\lambda_j}e^{i\beta_j n\lambda_j}
e^{-n\lambda_js}.
\end{equation}

Let \(h=-G'/G\). The Euler representation gives

\begin{equation}\label{eq:ob4-5}
-M'/M=-\tfrac12\log A+h+D_{\rm ar}.
\end{equation}

Along the positive real axis, \(R,h\to0\), while \(C,D_{\rm ar}\) are exponentially small. Hence \(a=\tfrac12\log A\) and

\begin{equation}\label{eq:ob4-6}
R(s)-h(s)=D_{\rm ar}(s)-C(s)=O(e^{-cs}).
\end{equation}

The gamma recurrence gives

\begin{equation}\label{eq:ob4-7}
h(s+2)-h(s)=-\frac d s+\frac d{s+1}.
\end{equation}

Consequently the rational function

\begin{equation}\label{eq:ob4-8}
R(s+2)-R(s)+\frac d s-\frac d{s+1}
\end{equation}

is exponentially small on the positive real axis and is therefore identically zero. Its possible poles from \(R(s)\) lie in \(0<\Re s<1\), those from \(R(s+2)\) lie in \(-2<\Re s<-1\), and the two gamma poles lie on the separate lines \(\Re s=0,-1\). Thus \(d=0\), and \(R(s+2)=R(s)\). A rational periodic function that tends to zero at infinity is zero. This leaves

\begin{equation}\label{eq:ob4-9}
D_{\rm ar}=C.
\end{equation}

Partition the finitely many \(\lambda_j\) according to positive rational commensurability. In each class choose a common \(\lambda>0\) such that \(\lambda_j=m_j\lambda\) with positive integers \(m_j\). Different classes never share a nonzero frequency. Within one class the coefficient at \(n\lambda\) has the form

\begin{equation}\label{eq:ob4-10}
\sum_{\sigma\in\Sigma}
e^{\sigma n\lambda}u_\sigma(n),
\end{equation}

where \(\Sigma\subset(0,1)\) is finite and each \(u_\sigma\) is a finite sum of unit-modulus exponentials. Indeed, the restriction \(m_j\mid n\) is a finite average of roots of unity to the \(n\)-th power.

By uniqueness of the generalized Dirichlet expansions, \eqref{eq:ob4-10} equals the coefficient in \(D_{\rm ar}\). That coefficient is uniformly bounded along the progression, because the progression meets powers of at most one prime, as in Proposition~\ref{prop:ob-simple-peeling}. Divide by \(e^{\sigma_{\max}n\lambda}\). All lower exponents tend to zero, and the bounded left side does too; hence \(u_{\sigma_{\max}}(n)\to0\). A finite exponential sum

\begin{equation}\label{eq:ob4-11}
u(n)=\sum_{\ell=1}^Lz_\ell e^{in\theta_\ell}
\end{equation}

with distinct phases modulo \(2\pi\) cannot tend to zero unless every coefficient vanishes: its Cesàro mean square tends to \(\sum|z_\ell|^2\). Remove this largest exponent and repeat through the finite set \(\Sigma\). Every coefficient in every class is zero. Thus \(C=D_{\rm ar}=0\).

Now \eqref{eq:ob4-5} gives \(M(s)=e^{as+b}\). The reflection equation forces \(a=0\), so \(M\) is constant. At a prime in \(\Omega\), \eqref{eq:Dar} has coefficient \(a(p,1)=4\), which is incompatible with \(D_{\rm ar}=0\). \end{proof}

This proposition gives a precise strengthened target for any off-central argument: proving a finite progression description of the \emph{whole} divisor would suffice, even if the progressions lie on several different vertical lines. Boundary positive definiteness alone does not supply that target.

\subsection{Exact prime Euler factors with a boundary lattice}

The following construction preserves much more than an abstract positive Laplace transform. It uses rational primes, exactly the local factors \eqref{eq:Q}, bounded prime-power traces, and ordinary Dirichlet coefficients satisfying a Ramanujan bound. It has a meromorphic continuation through its outer line, whose entire boundary divisor is a single positive lattice comb. It shows concretely why those facts alone do not identify the Euler product with its boundary factor.

\begin{proposition}[An exact prime boundary example]\label{prop:ob-primes}

For every positive integer \(m\), there is an infinite set \(\Omega\) of odd primes such that \(Q_\Omega\) has all of the following properties:

\begin{enumerate}
\def\labelenumi{\arabic{enumi}.}

\item
  Its abscissa of absolute convergence is \(\alpha=3/5\), and \(\sum_{p\in\Omega}p^{-1}<\infty\).
\item
  Its local factors are exactly \(((1+p^{-s})/(1-p^{-s}))^2\), so its normalized logarithmic prime-power coefficients are \(4\) for odd powers and zero for even powers.
\item
  Its ordinary coefficients are nonnegative integers and satisfy \(a(n)\le d_4(n)\ll_\epsilon n^\epsilon\).
\item
  It continues meromorphically to \(\Re s>4/15\), where

\begin{equation}\label{eq:ob5-1}
Q_\Omega(s)=(1-2^{3/5-s})^{-m}e^{G(s)}
\end{equation}

for a holomorphic function \(G\).
\item In that half-plane it has no zeros and has precisely the poles

\begin{equation}\label{eq:ob5-2}
s=\frac35+\frac{2\pi i k}{\log2},\qquad k\in\mathbb Z,
\end{equation}

each of order \(m\).
\item The zero-free factor after removing the displayed boundary factor has logarithmic Dirichlet coefficients

\begin{equation}\label{eq:ob5-3}
-m(\log2)2^{3n/5}
\end{equation}

at every frequency \(n\log2\), and hence does not retain positivity.
\end{enumerate}

\end{proposition}

\begin{proof}\textit{Published input.} \citet[p.~562, final displayed estimate]{bhp2001} prove

\begin{equation}\label{eq:ob5-4}
\pi(x+x^{21/40})-\pi(x)
\ge\frac9{100}\frac{x^{21/40}}{\log x}
\end{equation}

for all sufficiently large \(x\). Only the weak consequence

\begin{equation}\label{eq:ob5-5}
\pi(X+X^{2/3})-\pi(X)\gg X^{2/3}/\log X
\end{equation}

is needed. To see it directly, pack disjoint intervals starting at \(X+j(2X)^{21/40}\); each contains its shorter interval of length \(x^{21/40}\), and there are \(\gg X^{2/3-21/40}\) of them within \([X,X+X^{2/3}]\). Multiplying \eqref{eq:ob5-4} by this number proves \eqref{eq:ob5-5}. The constant \(9/100\) is not needed elsewhere.

\textit{Construction.} Set

\begin{equation}\label{eq:ob5-6}
X_n=2^n,\quad \lambda=\log2,\quad
\alpha=\frac35,\quad\beta=\frac23,
\quad T_n=\frac{m\lambda}{4}X_n^\alpha.
\end{equation}

For all sufficiently large \(n\), the interval

\begin{equation}\label{eq:ob5-7}
I_n=[X_n,X_n+X_n^\beta]
\end{equation}

contains primes of total logarithmic weight \(\gg X_n^\beta>T_n\). The intervals are disjoint. In increasing order, select primes from \(I_n\) until their logarithmic weights first reach \(T_n\). Let \(\Omega_n\) be this finite selection and \(\Omega\) their union, omitting all finitely many initial intervals. All chosen primes are odd. The overshoot is at most the logarithm of one prime, so

\begin{equation}\label{eq:ob5-8}
W_n:=\sum_{p\in\Omega_n}\log p
=T_n+O(\log X_n).
\end{equation}

In particular \(W_n\asymp X_n^\alpha\), and \(\#\Omega_n\asymp X_n^\alpha/\log X_n\). These estimates prove convergence of the prime sum for \(\Re s>\alpha\), divergence at \(\alpha\), and convergence at \(s=1\). The positive Euler product therefore has abscissa exactly \(\alpha\), as does its logarithmic derivative. At \(\alpha\), even \(\sum_{p\in\Omega}p^{-\alpha}\) diverges by comparison with \(\sum_n1/n\).

For any compact set of \(s\)'s, the mean value estimate along the real segment from \(X_n\) to \(p\in I_n\) gives

\begin{equation}\label{eq:ob5-9}
|p^{-s}-X_n^{-s}|
\ll_K X_n^{-\Re s-1}(p-X_n).
\end{equation}

Consequently

\begin{equation}\label{eq:ob5-10}
\sum_{p\in\Omega_n}(\log p)p^{-s}
=T_nX_n^{-s}
+O_K\bigl(X_n^{\alpha+\beta-1-\Re s}
+(\log X_n)X_n^{-\Re s}\bigr).
\end{equation}

The sum of errors converges normally when

\begin{equation}\label{eq:ob5-11}
\Re s>\alpha+\beta-1=\frac4{15}.
\end{equation}

The odd prime powers of exponent at least three converge normally for \(\Re s>\alpha/3=1/5\), since they are bounded on compact sets by a constant times \(\sum_{p\in\Omega}(\log p)p^{-3\Re s}\). Including the finitely many omitted geometric terms, \eqref{eq:ob5-10} therefore yields

\begin{equation}\label{eq:ob5-12}
D_\Omega(s)
=\frac{m\lambda}{e^{\lambda(s-\alpha)}-1}+E(s),
\end{equation}

with \(E\) holomorphic in \(\Re s>4/15\). Initially this identity is an identity of absolutely convergent series in \(\Re s>\alpha\), and its right side supplies the continuation.

Choose a holomorphic primitive \(G\) of \(-E\) on that simply connected half-plane. Since

\begin{equation}\label{eq:ob5-13}
-\frac d{ds}\log(1-e^{-\lambda(s-\alpha)})^{-m}
=\frac{m\lambda}{e^{\lambda(s-\alpha)}-1},
\end{equation}

the quotient of \(Q_\Omega\) and the right side of \eqref{eq:ob5-1} is constant in the original half-plane. Choose the additive constant of \(G\) to make this constant one. This proves \eqref{eq:ob5-1}--\eqref{eq:ob5-2}, including all divisor orders.

The local expansion is

\begin{equation}\label{eq:ob5-14}
\left(\frac{1+x}{1-x}\right)^2
=1+4\sum_{k\ge1}k x^k.
\end{equation}

Since \(4k\le\binom{k+3}{3}=d_4(p^k)\), multiplicativity gives the claimed ordinary coefficient bound. Finally, \(2\notin\Omega\), so the original logarithmic derivative has no term at any \(n\log2\). Subtracting the model derivative in \eqref{eq:ob5-12} gives exactly \eqref{eq:ob5-3}. \end{proof}

\paragraph{Scope of the example.}

There is no claim that \(\Omega\) comes from an \(S_5\) representation, and no claim of global meromorphic continuation or order at most one. In fact the reflection functional equation required in \eqref{eq:FE} \emph{fails already on the displayed continuation domain}: the real pole at \(3/5\) would require a pole at \(2/5\), whereas \eqref{eq:ob5-1} is holomorphic and nonzero at \(2/5>4/15\). Multiplication by a nonvanishing conductor exponential cannot repair that missing reflected divisor.

Thus Proposition~\ref{prop:ob-primes} is a counterexample to a boundary-only extension, not to the complete arithmetic conjecture. It establishes that exact prime Euler factors, the Ramanujan coefficient bound, integral boundary orders, Cohen's lattice conclusion, and continuation a fixed distance across the boundary do not by themselves force a finite exceptional set.

\subsection{A symmetric order-one example without an interior gap}

Local finiteness of a divisor gives a horizontal gap on every bounded-height compact set after one removes its boundary points. It does not give a gap uniform in height, even for a globally meromorphic function of order one with the exact reflection symmetry.

\begin{proposition}[No uniform interior gap]\label{prop:ob-gap}

There is a real-type meromorphic function \(M_*\) of order at most one such that \(M_*(1-s)=M_*(s)\), all its zeros and poles lie in \(2/5\le\Re s\le3/5\), and its right boundary profile is \(\mathbf1_{\{0\}}\), while both zeros and poles strictly inside the strip approach the line \(\Re s=3/5\) at unbounded heights.

\end{proposition}

\begin{proof} Set \(z=s-1/2\), \(a=1/10\), and, for \(n\ge1\),

\begin{equation}\label{eq:ob6-1}
u_n=a-\frac1{n+12}+in,\qquad
v_n=a-\frac2{n+24}+i(n+1/2).
\end{equation}

Their real parts lie strictly between zero and \(a\). Define

\begin{equation}\label{eq:ob6-2}
P_u(z)=\prod_{n\ge1}
\left(1-\frac{z^2}{u_n^2}\right)
\left(1-\frac{z^2}{\overline u_n^2}\right),
\end{equation}

and define \(P_v\) in the same way. These products converge uniformly on compact sets since \(\sum|u_n|^{-2}\) and \(\sum|v_n|^{-2}\) converge. Moreover

\begin{equation}\label{eq:ob6-3}
|P_u(z)|, |P_v(z)|
\le\prod_{n\ge1}(1+|z|^2/n^2)^2
=\left(\frac{\sinh(\pi|z|)}{\pi|z|}\right)^2.
\end{equation}

Thus they are entire of exponential type and order at most one. Put

\begin{equation}\label{eq:ob6-4}
M_*(s)=
\frac{P_v(s-1/2)}{
((s-1/2)^2-a^2)P_u(s-1/2)}.
\end{equation}

Its numerator and denominator are entire of order at most one, so it is meromorphic of order at most one. Both products are even and of real type, giving the required symmetries. All listed zeros and poles are distinct: the nonreal poles have integral imaginary parts, and the nonreal zeros have half-integral imaginary parts. There are simple real poles at \(s=3/5,2/5\). The right-hand interior poles have coordinates

\begin{equation}\label{eq:ob6-5}
s=\frac35-\frac1{n+12}+in,
\end{equation}

and the corresponding zeros have coordinates

\begin{equation}\label{eq:ob6-6}
s=\frac35-\frac2{n+24}+i(n+1/2).
\end{equation}

Their conjugates and reflected points complete the divisor. This proves every asserted property. For real \(\sigma>3/5\), the conjugate factor pairs are positive and \((\sigma-1/2)^2-a^2>0\), so \(M_*(\sigma)>0\) as well. \end{proof}

No Dirichlet expansion or Euler product is claimed for \(M_*\). Its role is exact: the reflection equation, finite order, local finiteness, a positive simple real pole, and the permitted boundary profile do \emph{not} imply a uniform distance between the remaining divisor and the outer line.

\subsection{A reflected positive Dirichlet example without bounded traces}

For a complementary counterexample, let \(\chi_5\) be the even primitive quadratic character modulo five and put

\begin{equation}\label{eq:ob7-1}
w=3s-1,\qquad
Q_0(s)=\frac{\zeta(w)}{L(w,\chi_5)},\qquad
M_0(s)=5^{-w/2}Q_0(s).
\end{equation}

The completed degree-one functions have the same gamma factor and root number \(+1\). Since \(w(1-s)=1-w(s)\), their completed ratio satisfies

\begin{equation}\label{eq:ob7-2}
M_0(1-s)=M_0(s).
\end{equation}

The standard zeta and primitive Dirichlet functional equations used here are recorded in \citet[\S\S25.4 and 25.15, equations 25.4.1 and 25.15.5]{dlmf}. For this particular character, its Gauss sum is

\begin{equation}\label{eq:ob7-3}
2\cos(2\pi/5)-2\cos(4\pi/5)=\sqrt5,
\end{equation}

which checks the root number directly. The ratio is globally meromorphic of order at most one.

The Euler factors of \(Q_0\) are

\begin{equation}\label{eq:ob7-4}
\begin{cases}
(1-5^{1-3s})^{-1},&p=5,\\
(1+p^{1-3s})/(1-p^{1-3s}),&\chi_5(p)=-1,\\
1,&\chi_5(p)=1.
\end{cases}
\end{equation}

They give an ordinary Dirichlet series with nonnegative integer coefficients, supported on cubes, and a logarithmic derivative with nonnegative coefficients. It converges at \(s=1\), has abscissa \(2/3\), and has a simple real pole there because \(\zeta(w)\) has its simple pole at \(w=1\) while \(L(1,\chi_5)>0\). The last inequality follows directly by grouping the convergent series into blocks:
\[
L(1,\chi_5)=\sum_{n\ge0}
\left(\frac1{(5n+1)(5n+2)}
-\frac1{(5n+3)(5n+4)}\right)>0.
\]
Thus positivity, global order one, reflection, and convergence at one alone do not rule out an off-central real pole.

The crucial arithmetic bound fails explicitly. At an inert prime \(p\), the coefficient of \(p^{-3ks}\) in \(-Q_0'/Q_0\), for odd \(k\), is

\begin{equation}\label{eq:ob7-5}
6p^k\log p.
\end{equation}

The normalized prime-power traces are unbounded. For the fixed inert prime \(p=2\), the ordinary coefficient at \(2^{3k}\) is \(2^{k+1}\). Thus the Ramanujan estimate \(a(n)\ll_\epsilon n^\epsilon\) fails for every fixed \(\epsilon<1/3\), without needing any assertion about the distribution of inert primes. This example is therefore excluded by \eqref{eq:bound}. The unscaled ratio \(\zeta(s)/L(s,\chi_5)\), already discussed in Proposition~\ref{prop:counterexample}, has its pole at one; the substitution here moves the abscissa below one while preserving reflection about \(1/2\), at exactly the cost displayed in \eqref{eq:ob7-5}.

\paragraph{Joint scope.}
No example in this appendix satisfies the complete arithmetic comparison package while violating the conclusion of Theorem~\ref{thm:main}. Proposition~\ref{prop:ob-primes} preserves the exact prime Euler factors and bounded traces but fails reflection. Proposition~\ref{prop:ob-gap} preserves reflection and order one but asserts no Euler structure. The scaled classical ratio has reflection and positivity but fails the local coefficient bound. The unresolved problem is to use the global functional equation and bounded prime Euler structure jointly to control the interior divisor. The obstruction to canonical positive removal is not an obstruction to every possible analytic reduction.


\section{Tensor identities and the scope of further transfers}\label{sec:functorial-identities}

This appendix records exact identities which distinguish the exterior-square argument used in the main theorem from stronger functorial transfers that would settle the rank-five comparison directly. Write $V=\rho_4$, $R=\rh$, and $W=\rho_6$.

\begin{proposition}[Finite identities]\label{prop:finite-tensors}
The following identities hold as complex $S_5$-representations:
\begin{align*}
R\otimes R&=\one\oplus V\oplus V\eta\oplus R\oplus R\eta\oplus W,\\
V\otimes R&=V\oplus R\oplus R\eta\oplus W,\\
W\otimes R&=V\oplus V\eta\oplus R\oplus R\eta\oplus W^{\oplus2},\\
\operatorname{Sym}^2V&=\one\oplus V\oplus R,\\
\operatorname{Sym}^2R&=\one\oplus V\oplus R\oplus R\eta,\qquad
\wedge^2R=V\eta\oplus W.
\end{align*}
\end{proposition}

\begin{proof}
For a permutation $g$, let $a$ and $b$ be its numbers of fixed letters and two-cycles. The characters of $V$ and $R$ are $a-1$ and $a(a-3)/2+b$, respectively: the permutation character on unordered pairs is $\binom a2+b$, and subtracting the five-letter permutation character gives $R$. The character of $W$ is
\[
\chi_W(g)=\frac{\chi_V(g)^2-\chi_V(g^2)}2.
\]
Use the class order in \eqref{eq:rhochar}, with sizes $1,10,15,20,30,24,20$. Multiplying characters gives the first three rows; using
\[
\chi_{\operatorname{Sym}^2 X}(g)=\frac{\chi_X(g)^2+\chi_X(g^2)}2,
\qquad
\chi_{\wedge^2 X}(g)=\frac{\chi_X(g)^2-\chi_X(g^2)}2
\]
gives the last two. Direct substitution verifies equality on every class, which proves equality of the semisimple representations. For example the symmetric square of $V$ has character $(10,4,2,1,0,0,1)$, the character of the permutation action on unordered pairs. Thus it is exactly $\one\oplus V\oplus R$.
\end{proof}

\begin{proposition}[Exceptional tensor comparisons]\label{prop:tensor-cancellation}
At an exceptional prime put $T=p^{-s}$ and $\Phi_p=((1+T)/(1-T))^2$. The local quotients of the self tensor, symmetric square, and exterior square of the automorphic rank-five parameter by their Artin counterparts are respectively
\[
\Phi_p^2,\qquad\Phi_p,\qquad\Phi_p.
\]
Tensoring that parameter with either $V$ or $W$ instead gives quotient one. Twisting the original rank-five comparison by $\eta$ gives $\Phi_p^{-1}$.
\end{proposition}

\begin{proof}
At a four-cycle the spectra of $V$ and $W$ are
\[
\{1,i,-1,-i\},\qquad\{1,-1,i,i,-i,-i\}.
\]
The two rank-five spectra are $r=\{1,-1,-1,i,-i\}$ and $a=\{1,1,1,i,-i\}$. Counting products, with coordinates ordered by the eigenvalues $1,i,-1,-i$, gives
\[
\begin{array}{c|rrrr|rrrr}
 &\multicolumn{4}{c|}{\text{Artin}}&\multicolumn{4}{c}{\text{automorphic}}\\
\text{self tensor}&7&6&6&6&11&6&2&6\\
\operatorname{Sym}^2&5&3&4&3&7&3&2&3\\
\wedge^2&2&3&2&3&4&3&0&3\\
V\text{ tensor}&5&5&5&5&5&5&5&5\\
W\text{ tensor}&7&8&7&8&7&8&7&8.
\end{array}
\]
Taking reciprocal Euler polynomials proves each assertion. Finally $\eta(p)=-1$, so twisting replaces $T$ by $-T$ and inverts $\Phi_p$.
\end{proof}

In particular the two rank-five terms in the Artin decomposition of $R\otimes R$ have comparison quotients $\Phi_p$ and $\Phi_p^{-1}$, which cancel. The remaining self Rankin--Selberg quotient is $\Phi_p^2$; the holomorphy of its numerator does not exclude poles of the quotient at zeros of its denominator. By contrast, $\wedge^2R$ contains neither rank-five term, and its intrinsically entire Hecke function is the separate function $C$ used in Proposition~\ref{prop:commonpole}.

\begin{proposition}[Two conditional transfer equivalences]\label{prop:transfer-equivalence}
Let $\pi_4$ be the matched cuspidal representation for $V$, and let $\tau$ be the matched cuspidal representation over $E$ for either faithful three-dimensional $A_5$-representation $U=\operatorname{Ad}\varrho$. Each of the following is equivalent to the existence of a unitary cuspidal weak automorphic realization of $R$:
\begin{enumerate}[label=(\roman*)]
\item a unitary isobaric weak symmetric-square transfer of $\pi_4$ to $\GL_{10}/\Q$;
\item a unitary isobaric weak $\operatorname{Asai}^+$ transfer of $\tau$ to $\GL_9/\Q$.
\end{enumerate}
Here weak means equality of unramified local parameters outside a finite set. The plus convention for tensor induction is the one for which an odd element $g$ has trace $\chi_U(g^2)$.
\end{proposition}

\begin{proof}
The Rankin--Selberg inputs below are the unequal-rank boundary regularity of \citet[Proposition~3.3, p.~795]{jacquetshalika1981ii}, their equal-rank simple-pole criterion \citep[Proposition~3.6, pp.~802--804]{jacquetshalika1981ii}, and nonvanishing on the line one \citep[Theorem~3.7, p.~804]{jacquetshalika1981ii}. The last input is Shahidi's theorem, also stated directly in \citet[p.~462]{shahidi1980}. These results concern partial Rankin--Selberg functions of unitary cuspidal representations; applying them constituent by constituent gives the stated pole multiplicities for unitary isobaric sums.

The finite identity for (i) is $\operatorname{Sym}^2V=\one\oplus V\oplus R$. For (ii), the two three-dimensional characters of $A_5$ have values $3,-1,0,(1+\sqrt5)/2,(1-\sqrt5)/2$, with the last two interchanged. Their product is the character of $V|_{A_5}\oplus R|_{A_5}$. At odd elements of types $2,4,3\cdot2$, the plus tensor-induction traces are respectively $3,-1,0$, because their squares have types $1,2^2,3$. These are exactly the traces of $V\oplus R$. Therefore
\[
\operatorname{Asai}_{A_5}^{S_5,+}U=V\oplus R,
\qquad
\operatorname{Asai}_{A_5}^{S_5,-}U=V\eta\oplus R\eta.
\]

Suppose first that the transfer in (i) exists, and call it $\Sigma$. All of its good Satake eigenvalues have absolute value one, since they are the stated finite-image eigenvalues. In the Euler logarithm of $L^S(s,\Sigma\times\pi_4^\vee)$, higher prime powers are therefore bounded as real $s\downarrow1$. Chebotarev gives prime-term coefficient
\[
\langle\chi_{\one+V+R},\chi_V\rangle_{S_5}=1
\]
in front of $\log(1/(s-1))$. Rankin--Selberg pole theory and nonvanishing on the line one show that $\pi_4$ occurs with multiplicity one in the unitary isobaric decomposition of $\Sigma$. The same argument with the trivial character shows that $\one$ occurs once. Removing those two constituents leaves a unitary isobaric representation $\Sigma_5$ of rank five whose good parameters are exactly those of $R$.

If $\Sigma_5=\boxplus_j\sigma_j^{\boxplus m_j}$ with distinct unitary cuspidal $\sigma_j$, the pole order of its self Rankin--Selberg function is $\sum_jm_j^2$. The finite character computation instead gives
\[
\langle\chi_R,\chi_R\rangle_{S_5}=1.
\]
Thus there is exactly one constituent of multiplicity one, proving cuspidality. This proves the forward implication in (i). In (ii), the same argument extracts $\pi_4$ once from the rank-nine transfer and proves that its rank-five complement is cuspidal. Conversely a weak cuspidal realization $\sigma_5$ of $R$ gives the respective transfers
\[
\one\boxplus\pi_4\boxplus\sigma_5,
\qquad
\pi_4\boxplus\sigma_5.
\]
These have exactly the required good parameters by the finite identities.

A weak cuspidal realization $\sigma_5$ also settles the specified comparison with $\pi$. Their good parameters differ only at a finite set and at $\Omega$, and $\sum_{p\in\Omega}p^{-1}<\infty$. Their Euler logarithms therefore show that
\[
\log L^S(s,\pi\times\sigma_5^\vee)
-\log L^S(s,\sigma_5\times\sigma_5^\vee)
\]
remains bounded as real $s\downarrow1$; all higher prime powers are uniformly summable. The second function has a simple pole, so the first does as well. Rankin--Selberg pole theory gives $\pi\simeq\sigma_5$. The good discrepancy set is then finite, and Proposition~\ref{prop:package}(d) gives equality of the standard local factors at every place.
\end{proof}

These equivalences specify an additional transfer that would solve the target; they do not supply it. \citet[Proposition 3.1(b)]{asgariraghuram} explicitly distinguish the conjectural symmetric-square transfer from the established exterior-square transfer in the tensor-product setting. Over $E$, the particular tensor product of the two adjoint three-dimensional representations is already an isobaric sum of the known four- and five-dimensional $A_5$ constituents. Descending it with the correct inert-prime parameters is the same unresolved choice encountered in the rank-five construction. The analytic exterior-square holomorphy used in the main proof is available without making either additional transfer assumption.

\subsection{A symmetric-square bound subsumed by the Artin-factor criterion}

Symmetric-square holomorphy gives a further Dedekind-zeta zero at a comparison pole. The factorization below shows exactly why this does not strengthen the earlier Artin-factor variant in Remark~\ref{rem:weaker}.

\begin{proposition}[Symmetric-square condition and its redundancy]\label{prop:symmetric-square-redundancy}
Let $D_8=C_{S_5}((12)(34))$, let $T=L^{D_8}$, and put $D(s)=\zeta_T(s)$. Then
\[
\operatorname{Sym}^2R=\Ind_{D_8}^{S_5}\one,\qquad [T:\Q]=15,
\]
and $T$ has signature $(3,6)$ and infinite factor $\Gr(s)^9\Gr(s+1)^6$.
At every pole $z$ of $M$ of order $r>0$ in $0<\Rea z<1$,
\[
r\leq\min\{\ord_zB,\ord_zC,\ord_zD\}.
\]
Thus absence of a common off-central real zero of $B,C,D$ suffices for the all-place and entire-automorphy conclusions of Theorem~\ref{thm:main}, with no restriction on central multiplicities. However,
\begin{equation}\label{eq:symmetric-square-redundancy}
D(s)=L(s,V)L(s,R)\zeta_F(s),\qquad
\ord_zD\geq\ord_zL(s,R)\quad(0<\Rea z<1).
\end{equation}
The $D$ pole bound is consequently already implied by the prior $L(s,R)$ bound; the three-function condition is subsumed by Remark~\ref{rem:weaker}.
\end{proposition}

\begin{proof}
The centralizer of $(12)(34)$ fixes 5, independently interchanges the two letters in each of its two orbits, and may interchange the orbits. It has order $2^2\cdot2=8$ and is dihedral, generated by $(1324)$ and $(12)$. Its cosets identify equivariantly with the fifteen double transpositions, by $xD_8\mapsto x(12)(34)x^{-1}$. A fixed point of $g$ in this action is exactly a double transposition commuting with $g$.

The symmetric-square character formula in Proposition~\ref{prop:finite-tensors} gives the full calculation
\[
\begin{array}{c|rrrrrrr}
\text{cycle type}&1&2&2^2&3&4&5&3\cdot2\\ \hline
\chi_R(g^2)&5&5&5&-1&1&0&-1\\
\chi_{\operatorname{Sym}^2R}(g)&15&3&3&0&1&0&0\\
\text{fixed double transpositions}&15&3&3&0&1&0&0.
\end{array}
\]
For completeness, $(12)$ commutes with exactly $(12)(34),(12)(35),(12)(45)$. The centralizer of $(12)(34)$ contains exactly the three double transpositions on $\{1,2,3,4\}$. A three-cycle or a permutation of type $3\cdot2$ has centralizer $C_3\times C_2$, containing no double transposition. A four-cycle has cyclic centralizer of order four, whose square is its unique double transposition. A five-cycle has cyclic centralizer of order five and contributes none. These counts prove the last row and hence the representation identity on every class. Artin induction gives $L(s,\operatorname{Sym}^2R)=\zeta_T(s)$, including ramified factors. Complex conjugation fixes three cosets and exchanges the remaining twelve in six pairs, giving signature $(3,6)$. Finally $\Gc(s)=\Gr(s)\Gr(s+1)$ gives the stated infinite factor.

At an exceptional prime the two rank-five spectra in Proposition~\ref{prop:package} have symmetric-square multiplicities
\[
\begin{array}{c|rrrr}
\text{eigenvalue}&1&i&-1&-i\\ \hline
\text{Artin}&5&3&4&3\\
\text{automorphic}&7&3&2&3 .
\end{array}
\]
These count products with indices $i\leq j$, including diagonal products. With $x=p^{-s}$ the two local factors are
\[
D_p=(1-x)^{-5}(1+x)^{-4}(1+x^2)^{-3},\qquad
\mathcal H_p=(1-x)^{-7}(1+x)^{-2}(1+x^2)^{-3}.
\]
Therefore $\mathcal H_p/D_p=((1+x)/(1-x))^2$, with the same quotient direction as $M$. At every other undeleted prime the parameters match.

The external analytic input is \citet[Theorem~7.1, pp.~49--50]{takeda2015}: for a cuspidal representation with unitary central character and a unitary character $\chi$, the partial twisted symmetric-square function is holomorphic except for possible simple poles at 0 and 1, permitted only if $\chi^r\omega_\pi^2=1$. The deleted set can consist exactly of the infinite places, those over 2, and the places where $\pi$ or $\chi$ ramifies. Apply this with rank five and $\chi=1$. In particular it supplies holomorphy throughout the open strip. We use this analytic theorem as a published input; it neither asserts nonvanishing there nor constructs a $\GL_{15}$ transfer.

Take a finite set $S$ containing those places, all primes ramified in $L$, and the original finite exceptions. Enlarging Takeda's set deletes unramified factors by multiplication by entire polynomials $\prod_{i\leq j}(1-\alpha_{p,i}\alpha_{p,j}p^{-s})$, so
\[
\mathcal H^S(s)=L^S(s,\pi,\operatorname{Sym}^2)
\]
remains holomorphic in the strip. These added polynomials are nonzero for $\Rea s>0$, since $\pi$ is tempered. On the Artin side every local eigenvalue on inertia invariants is a root of unity, or the invariant space is absent. Thus $D^S/D$ is a finite product of polynomials $1-\lambda p^{-s}$ with $|\lambda|=1$, and is holomorphic and nowhere zero for $\Rea s>0$. The finitely removed factors of $Q_\Omega$ likewise have no zero or pole there, since $1\pm p^{-s}\ne0$.

Absolute convergence in $\Rea s>1$ and the local calculation give
\[
\mathcal H^S(s)=D^S(s)Q_{\Omega\setminus S}(s).
\]
The equality continues meromorphically. The established completed comparison and the preceding deletion factors therefore give, on the open strip,
\[
M(s)=u(s)\frac{\mathcal H^S(s)}{D(s)}
\]
with $u$ holomorphic and nowhere zero. At a pole of order $r$ this implies
\[
\ord_zD=r+\ord_z\mathcal H^S\geq r.
\]
Together with Proposition~\ref{prop:commonpole} this proves the three order bounds at every complex point in the strip. The pole of $\zeta_T$ at 1 is outside the domain. Excluding off-central real triple zeros excludes all such real poles of $M$, and the same central-divisor argument used in Theorem~\ref{thm:main} gives its conclusions.

To establish the precise limitation, Proposition~\ref{prop:finite-tensors} and \eqref{eq:factorization} give
\[
\zeta_T=\zeta_{\Q}L(V)L(R)L(R\eta)=L(V)L(R)\zeta_F.
\]
The functions $L(V)$ and $\zeta_F$ are holomorphic in the strip, proving \eqref{eq:symmetric-square-redundancy}. The existing standard comparison already gives $\ord_zL(s,R)=r+\ord_zL(s,\pi)\geq r$ at a pole of $M$. Thus it proves the $D$ bound without the symmetric-square analytic theorem.

More explicitly, write $Z(f)$ for the set of points of positive meromorphic order in the open strip. Since $B=L(\eta)L(R)$ with $L(\eta)$ entire, the exact inclusions are
\[
Z(L(R))\cap Z(C)
\ \subseteq\ Z(B)\cap Z(C)\cap Z(D)
\ \subseteq\ Z(B)\cap Z(C).
\]
They remain valid after restriction to off-central real points. Hence the triple-zero exclusion is weaker than the $B,C$-only formulation but at least as restrictive as the already available $L(R),C$ formulation. No new criterion beyond that Artin-factor variant follows.
\end{proof}

\begin{remark}[Formal and arithmetic scope]
The displayed identities and holomorphy constraints alone do not imply that every common $B,C$ zero is a $D$ zero. As a local-germ example, assign $L(\eta)=L(V\eta)=t$ and $L(R)=L(R\eta)=L(V)=L(W)=\zeta_{\Q}=\zeta_F=M=L(\pi)=L(\pi\eta)=\mathcal H^S=1$. Then $B=C=t$ and $D=1$, consistently with these identities. This is only a formal example, not an $S_5$ field or a global system of Euler products. Any arithmetic common $B,C$ zero missed by $D$ must be a zero of $L(\eta)$, by \eqref{eq:symmetric-square-redundancy}. No strict arithmetic zero-set inclusion or newly certified field is established.
\end{remark}


\section*{Method and public-domain dedication}

The mathematics and text were produced with GPT-6 Astra in ChatGPT Pro mode. The accompanying records report targeted checks of cited inputs by Claude Opus 5.5 with effort set to high. Those reported checks are used only for the inputs they address. The central analytic argument, the common-zero divisor bound, and the boundary-divisor results received separate reconstructions from their statements by other instances of GPT-6 Astra before comparison with the proposed proofs. The explicit numerical certificates were evaluated by separately written MPFR and Arb implementations, with different exact finite-field and trace calculations. These are checks by language models; no human review was performed.

The original text is dedicated to the public domain under \href{https://creativecommons.org/publicdomain/zero/1.0/}{CC0 1.0}. Reproducible material and verification records are available through the \href{\workbenchurl}{zeta and $L$-functions workbench} and the \href{\recordurl}{mathematics workbench records}.

\begin{thebibliography}{99}
\bibitem[Jacquet and Shalika(1981)]{jacquetshalika1981ii}
Jacquet, H., and Shalika, J. A. (1981).
On Euler products and the classification of automorphic forms II.
\emph{American Journal of Mathematics}, \textbf{103}(4), 777--815.
Proposition 3.3, p.~795; Proposition 3.6, pp.~802--804; Theorem 3.7, p.~804.

\bibitem[Shahidi(1980)]{shahidi1980}
Shahidi, F. (1980). On nonvanishing of $L$-functions.
\emph{Bulletin of the American Mathematical Society (New Series)}, \textbf{2}(3), 462--464.
Theorem, p.~462.

\bibitem[Takeda(2015)]{takeda2015}
Takeda, S. (2015).
\emph{On a certain metaplectic Eisenstein series and the twisted symmetric square $L$-function}.
arXiv:1506.04791v1 (15 June 2015), Theorem 7.1, pp.~49--50.
\url{https://arxiv.org/abs/1506.04791v1}.

\bibitem[Asgari and Raghuram(2007)]{asgariraghuram}
Asgari, M., and Raghuram, A. (2007).
\emph{A cuspidality criterion for the exterior square transfer of cusp forms on $\GL(4)$}.
\url{https://arxiv.org/abs/0712.4315}.

\bibitem[Belt(2012)]{belt}
Belt, D. (2012). \emph{On the holomorphy of exterior-square $L$-functions}.
Preprint of doctoral dissertation, arXiv:1108.2200v7.
\url{https://arxiv.org/abs/1108.2200v7}.

\bibitem[Kim(1999)]{kim1999}
Kim, H. H. (1999). Langlands--Shahidi method and poles of automorphic $L$-functions: Application to exterior square $L$-functions.
\emph{Canadian Journal of Mathematics}, \textbf{51}(4), 835--849.
\href{https://doi.org/10.4153/CJM-1999-036-0}{doi:10.4153/CJM-1999-036-0}.

\bibitem[Pilloni and Stroh(2016)]{pillonistroh}
Pilloni, V., and Stroh, B. (2016).
Surconvergence, ramification et modularit\'e.
\emph{Ast\'erisque}, \textbf{382}, 195--266.
\href{https://doi.org/10.24033/ast.1004}{doi:10.24033/ast.1004}.

\bibitem[Khare, La Rosa, and Wiese(2023)]{kharelarosawiese}
Khare, C. B., La Rosa, A. F., and Wiese, G. (2023).
Splitting fields of $X^n-X-1$ (particularly for $n=5$), prime decomposition and modular forms.
\emph{Expositiones Mathematicae}, \textbf{41}(3), 475--491.
\href{https://doi.org/10.1016/j.exmath.2023.02.007}{doi:10.1016/j.exmath.2023.02.007}.

\bibitem[Baker et al.(2001)]{bhp2001}
Baker, R. C., Harman, G., \& Pintz, J. (2001). The difference between consecutive primes, II.
\emph{Proceedings of the London Mathematical Society, 83}(3), 532--562.
\url{https://doi.org/10.1112/plms/83.3.532}.

\bibitem[Booker(2006)]{booker2006}
Booker, A. R. (2006). Artin's conjecture, Turing's method, and the Riemann hypothesis.
\emph{Experimental Mathematics, 15}(4), 385--407.
\url{https://doi.org/10.1080/10586458.2006.10128976}.

\bibitem[Booker(2015)]{booker2015}
Booker, A. R. (2015). $L$-functions as distributions.
\emph{Mathematische Annalen, 363}(1--2), 423--454.
\url{https://doi.org/10.1007/s00208-015-1178-z}.
Author manuscript: \url{https://arxiv.org/abs/1308.3067}.

\bibitem[Calegari(2013)]{calegari}
Calegari, F. (2013). The Artin conjecture for some $S_5$-extensions.
\emph{Mathematische Annalen, 356}, 191--207.
\url{https://doi.org/10.1007/s00208-012-0839-4}.
Author-hosted publisher text: \url{https://www.math.uchicago.edu/~fcale/papers/S5.pdf}.

\bibitem[Cohen(1960)]{cohen}
Cohen, P. J. (1960). On a conjecture of Littlewood and idempotent measures.
\emph{American Journal of Mathematics, 82}(2), 191--212.
\url{https://doi.org/10.2307/2372731}.

\bibitem[Dikranjan(2018)]{dikranjan2018}
Dikranjan, D. (2018). \emph{Introduction to topological groups} (Version of 26 February 2018). University of Udine lecture notes.
\url{https://users.dimi.uniud.it/~dikran.dikranjan/ITG.pdf}.

\bibitem[Dwyer(2014)]{dwyer2014}
Dwyer, J. (2014). Real zeros of Artin $L$-functions corresponding to five-dimensional $S_5$-representations.
\emph{Bulletin of the London Mathematical Society, 46}(1), 51--58.
\url{https://doi.org/10.1112/blms/bdt075}.

\bibitem[Green and Sanders(2008)]{greensanders}
Green, B., \& Sanders, T. (2008). A quantitative version of the idempotent theorem in harmonic analysis.
\emph{Annals of Mathematics, 168}(3), 1025--1054.
\url{https://doi.org/10.4007/annals.2008.168.1025}.

\bibitem[Kolountzakis(2000)]{kolountzakis}
Kolountzakis, M. N. (2000). On the structure of multiple translational tilings by polygonal regions.
\emph{Discrete \& Computational Geometry, 23}(4), 537--553.
\url{https://doi.org/10.1007/s004540010014}.
Author manuscript: \url{https://arxiv.org/abs/math/9904065}.

\bibitem[K{\"o}rner(2008)]{korner2008}
K{\"o}rner, T. W. (2008). \emph{Topological groups} (Part III lecture notes, 8 March 2008). University of Cambridge.
\url{https://www.dpmms.cam.ac.uk/~twk10/Topg.pdf}.

\bibitem[Liu(2026)]{liu2026}
Liu, B. (2026). \emph{Joint marked moments and optimized zero-free half-planes for unitary Hecke $L$-functions} (Version 2). Hexagon preprint.
\url{https://hexagonmath.org/2611.00421v2}.

\bibitem[Miller(2002)]{miller}
Miller, S. D. (2002). The highest-lowest zero and other applications of positivity.
\emph{Duke Mathematical Journal, 112}(1), 83--116.
\url{https://doi.org/10.1215/S0012-9074-02-11213-7}.
Author manuscript: \url{https://arxiv.org/abs/math/0112196}.

\bibitem[MPFR Project(2025)]{mpfr}
MPFR Project. (2025). \emph{GNU MPFR 4.2.2 manual}.
\url{https://www.mpfr.org/mpfr-current/mpfr.html}.

\bibitem[NIST(2026)]{dlmf}
National Institute of Standards and Technology. (2026).
\emph{NIST Digital Library of Mathematical Functions} (Version 1.2.8), \S\S25.4 and 25.15.
\url{https://dlmf.nist.gov/25.4}; \url{https://dlmf.nist.gov/25.15}.

\bibitem[Ramakrishnan(2002)]{ramakrishnan}
Ramakrishnan, D. (2002). Modularity of solvable Artin representations of $\mathrm{GO}(4)$-type.
\emph{International Mathematics Research Notices, 2002}(1), 1--54.
Author manuscript: \url{https://www.its.caltech.edu/~dinakar/papers/go1.pdf}.

\bibitem[Serre(1979)]{serrelocal}
Serre, J.-P. (1979). \emph{Local fields} (M. J. Greenberg, Trans.; Graduate Texts in Mathematics, Vol. 67).
Springer. \url{https://doi.org/10.1007/978-1-4757-5673-9}.

\end{thebibliography}
\end{document}
